\documentclass[11pt,a4paper]{article}
\usepackage[utf8]{inputenc}
\usepackage{amsmath,amssymb,amsthm}
\usepackage{algorithm,algorithmic}
\usepackage{booktabs,array,longtable}
\usepackage{hyperref}
\usepackage{geometry}
\usepackage{enumitem}
\usepackage{xspace}
\usepackage{mathtools}
\usepackage{multirow}
\usepackage{float}
\usepackage{caption}
\usepackage{subcaption}
\usepackage{tikz}
\usepackage{pgfplots}
\pgfplotsset{compat=1.18}
\usepackage{xcolor}
\usepackage{colortbl}
\usepackage{stmaryrd}
\usepackage{mathrsfs}

% Page geometry
\geometry{margin=2.5cm}

% Theorem environments
\newtheorem{theorem}{Theorem}[section]
\newtheorem{lemma}[theorem]{Lemma}
\newtheorem{definition}[theorem]{Definition}
\newtheorem{assumption}[theorem]{Assumption}
\newtheorem{remark}[theorem]{Remark}
\newtheorem{corollary}[theorem]{Corollary}
\newtheorem{claim}[theorem]{Claim}

% Custom commands
\newcommand{\polarKem}{\textsf{Polar-KEM}\xspace}
\newcommand{\CoreKEM}{\textsf{Core.KEM}\xspace}
\newcommand{\lipPrb}{\textsf{LIP}\xspace}
\newcommand{\BDD}{\textsf{BDD}\xspace}
\newcommand{\CVP}{\textsf{CVP}\xspace}
\newcommand{\usvp}{\textsf{uSVP}\xspace}
\newcommand{\BKZ}{\textsf{BKZ}\xspace}
\newcommand{\scCanc}{\textsf{SC}\xspace}
\newcommand{\scCancL}{\textsf{SCL}\xspace}
\newcommand{\FO}{\textsf{FO}\xspace}
\newcommand{\KDF}{\textsf{KDF}\xspace}
\newcommand{\Extract}{\textsf{Extract}\xspace}
\newcommand{\polarLIP}{\textsf{polar-LIP}\xspace}
\newcommand{\hidPrbl}{\textsf{HID}\xspace}
\newcommand{\INDCPA}{\textsf{IND-CPA}\xspace}
\newcommand{\INDCCA}{\textsf{IND-CCA2}\xspace}
\newcommand{\QROM}{\textsf{QROM}\xspace}
\newcommand{\DFP}{\textsf{DFP}\xspace}
\newcommand{\CTSelect}{\textsf{CTSelect}\xspace}
\newcommand{\CTEqual}{\textsf{CTEqual}\xspace}
\newcommand{\NIST}{\textsf{NIST}\xspace}
\newcommand{\MLKEM}{\textsf{ML-KEM}\xspace}
\newcommand{\hwkRef}{\textsf{HAWK}\xspace}
\newcommand{\FALCON}{\textsf{FALCON}\xspace}
\newcommand{\Kyber}{\textsf{Kyber}\xspace}
\newcommand{\BIKE}{\textsf{BIKE}\xspace}
\newcommand{\bit}{\{0,1\}}
\newcommand{\Z}{\mathbb{Z}}
\newcommand{\R}{\mathbb{R}}
\newcommand{\Q}{\mathbb{Q}}
\newcommand{\F}{\mathbb{F}}
\DeclareMathOperator{\GL}{GL}
\DeclareMathOperator{\sL}{SL}
\DeclareMathOperator{\negl}{negl}
\newcommand{\sample}{\xleftarrow{\$}}
\DeclareMathOperator{\Ext}{Ext}
\DeclareMathOperator{\Adv}{Adv}
\DeclareMathOperator{\Aut}{Aut}
\DeclareMathOperator{\lrnPrb}{LWE}
\DeclareMathOperator{\VPb}{SVP}
\newcommand{\pk}{\mathsf{pk}}
\newcommand{\sk}{\mathsf{sk}}
\newcommand{\ct}{\mathsf{ct}}
\DeclareMathOperator{\polarScDec}{PolarSCDecode}
\DeclarePairedDelimiter{\ceil}{\lceil}{\rceil}
\DeclarePairedDelimiter{\floor}{\lfloor}{\rfloor}
\DeclarePairedDelimiter{\norm}{\lVert}{\rVert}
\DeclarePairedDelimiter{\abs}{\lvert}{\rvert}

% Confidence tags
\newcommand{\heuTag}{\textcolor{orange}{\textbf{[Heuristic]}}}
\newcommand{\Conjectural}{\textcolor{red}{\textbf{[Conjectural]}}}
\newcommand{\prv}{\textcolor{green!60!black}{\textbf{[Proven]}}}

\begin{document}

%%%%%%%%%%%%%%%%%%%%%%%%%%%%%%%%%%%%%%%%%%%%%%%%%%%%%%%%%%%%%%%%%%%%%%%%%%%%%%%%
% TITLE PAGE
%%%%%%%%%%%%%%%%%%%%%%%%%%%%%%%%%%%%%%%%%%%%%%%%%%%%%%%%%%%%%%%%%%%%%%%%%%%%%%%%
\begin{titlepage}
\centering
\vspace*{2cm}

{\Huge\bfseries Polar-KEM\par}
\vspace{0.5cm}
{\LARGE\itshape A Lattice-Based Key Encapsulation Mechanism\\Using Polar-Code-Defined Lattices\par}

\vspace{2cm}

{\Large\bfseries Algorithm Specification\par}
\vspace{0.3cm}
{\large Version 1.0 -- Draft\par}

\vspace{1.5cm}

{\large Date: \today\par}

\vspace{3cm}

{\large
\textbf{Authors:}\\[0.5em]
\textsuperscript{1}Author Placeholder 1,
\textsuperscript{2}Author Placeholder 2,
\textsuperscript{1}Author Placeholder 3,\\
\textsuperscript{3}Author Placeholder 4,
\textsuperscript{2}Author Placeholder 5
\par}

\vspace{1cm}

{\normalsize
\textsuperscript{1}Institution Placeholder 1\\
\textsuperscript{2}Institution Placeholder 2\\
\textsuperscript{3}Institution Placeholder 3\\
\texttt{contact@polarkem.example.org}
\par}

\vfill

\begin{abstract}
\noindent
This document presents \polarKem, a post-quantum Key Encapsulation Mechanism (KEM) whose security is based on the Lattice Isomorphism Problem (LIP) over polar-code-defined lattices. Polar lattices, constructed via Construction~D from nested polar codes, admit efficient and reliable bounded-distance decoding using the successive cancellation (SC) decoder, enabling a clean separation between correct decryption and hard lattice problems. \polarKem builds on the blueprint of LIP-based cryptography but departs from prior constructions by using polar lattices instead of Barnes-Wall or orthogonal lattices, yielding improved decoding efficiency and concrete security estimates. The scheme is designed for IND-CCA2 security in the quantum random oracle model via a Fujisaki-Okamoto transform with implicit rejection. We propose three parameter sets---\textsf{PolarKEM-128}, \textsf{PolarKEM-256}, and \textsf{PolarKEM-512}---targeting NIST security levels 1, 3, and 5 respectively. Our analysis suggests that \polarKem achieves competitive key and ciphertext sizes while offering a novel security assumption based on the hardness of the polar-LIP problem.
\end{abstract}

\vfill
{\small\textit{This document is a submission to a cryptographic standards body.}}

\end{titlepage}

%%%%%%%%%%%%%%%%%%%%%%%%%%%%%%%%%%%%%%%%%%%%%%%%%%%%%%%%%%%%%%%%%%%%%%%%%%%%%%%%
% TABLE OF CONTENTS
%%%%%%%%%%%%%%%%%%%%%%%%%%%%%%%%%%%%%%%%%%%%%%%%%%%%%%%%%%%%%%%%%%%%%%%%%%%%%%%%
\tableofcontents
\newpage

%%%%%%%%%%%%%%%%%%%%%%%%%%%%%%%%%%%%%%%%%%%%%%%%%%%%%%%%%%%%%%%%%%%%%%%%%%%%%%%%
% INTRODUCTION
%%%%%%%%%%%%%%%%%%%%%%%%%%%%%%%%%%%%%%%%%%%%%%%%%%%%%%%%%%%%%%%%%%%%%%%%%%%%%%%%
\section{Introduction}\label{sec:intro}

\subsection{Motivation}\label{sec:motivation}

The advent of large-scale quantum computing threatens the security of widely deployed public-key cryptosystems based on integer factorization and discrete logarithms. In response, the cryptographic community has developed post-quantum alternatives, with lattice-based schemes receiving particular attention due to their efficiency, versatility, and strong security foundations.

Most lattice-based KEMs submitted to the \NIST Post-Quantum Cryptography standardization process---including \MLKEM, \Kyber, \BIKE, and others---base their security on variants of the Learning With Errors (LWE) problem or the Shortest Vector Problem (SVP). While these assumptions are well-studied, the diversification of hard problems strengthens the overall post-quantum ecosystem. The Lattice Isomorphism Problem (\lipPrb) offers an alternative foundation: given two isometric lattices, find the isometry between them. \lipPrb has been used in signature schemes (notably \hwkRef) but has seen limited exploration for KEM constructions.

\subsection{Why LIP-based KEM?}\label{sec:why-lip}

A KEM construction from \lipPrb proceeds by using the secret isometry as the private key and a public representative of the isometry class as the public key. Encryption embeds a short error vector into the public lattice; decryption uses the secret isometry to map the problem to an ``easy'' lattice where the error can be removed via efficient decoding.

The primary advantages of \lipPrb-based KEMs include:
\begin{enumerate}[leftmargin=2em]
\item \textbf{Unique security foundation:} Security rests on the hardness of finding isometries between lattices, a problem with algebraic structure distinct from LWE or SVP.
\item \textbf{Efficient decryption:} When the secret lattice admits fast decoding (e.g., via Construction~D from polar codes), decryption is substantially faster than generic lattice decoding.
\item \textbf{Tight security reductions:} Under appropriate assumptions, \lipPrb-based schemes admit tight reductions to well-defined lattice problems.
\end{enumerate}

\subsection{Why Polar Lattices?}\label{sec:why-polar}

Prior \lipPrb-based constructions have primarily used Barnes-Wall lattices or simple orthogonal lattices as the underlying ``easy'' lattice. Polar lattices offer several compelling advantages:

\begin{enumerate}[leftmargin=2em]
\item \textbf{Efficient decoding:} Polar lattices support successive cancellation (SC) decoding in $\mathcal{O}(N \log N)$ time, matching the efficiency of Barnes-Wall decoders while offering more flexible parameter choices.
\item \textbf{Flexible dimension and rate:} Unlike Barnes-Wall lattices (which are restricted to dimensions $N = 2^m$ with limited minimum-distance properties), polar lattices can be instantiated at any dimension and with fine-grained control over the code rate and minimum distance.
\item \textbf{Structured minimum distance:} The minimum distance of a polar lattice is precisely controlled by the polar code construction, yielding reliable bounded-distance decoding with analytically determined decoding radii.
\item \textbf{Parallelism:} The SC decoder for polar lattices is highly parallelizable, making \polarKem attractive for hardware implementations.
\end{enumerate}

\subsection{Relation to HAWK}\label{sec:relation-hawk}

\hwkRef\cite{HAWK} is a digital \emph{signature} scheme based on the Lattice Isomorphism Problem over orthogonal lattices. We emphasize that \polarKem is \textbf{not} related to \hwkRef as a KEM analogue; rather, both schemes share the broad methodology of \lipPrb-based cryptography:

\begin{itemize}[leftmargin=2em]
\item \hwkRef uses the \lipPrb for \emph{signatures}: the signer proves knowledge of a short vector in the secret lattice.
\item \polarKem uses the \lipPrb for \emph{key encapsulation}: the encapsulator embeds a secret into the public lattice, and the decapsulator recovers it using the secret isometry.
\end{itemize}

The algorithmic structures of \hwkRef (sign/verify) and \polarKem (encaps/decaps) are entirely different. In particular, \polarKem does \emph{not} use \hwkRef's Sign or Verify algorithms, does not require Gaussian sampling, and targets a completely different security goal.

\subsection{Relation to Barnes-Wall LIP-KEM Baseline}\label{sec:relation-bw}

A natural baseline for \polarKem is a hypothetical KEM built from \lipPrb over Barnes-Wall (BW) lattices. Both constructions use Construction~D with nested codes, but the code families differ:

\begin{itemize}[leftmargin=2em]
\item BW lattices use Reed-Muller codes, which have fixed rate progression and limited flexibility.
\item Polar lattices use polar codes, which offer near-arbitrary rate control and can be optimized for the channel at each level.
\end{itemize}

This flexibility allows \polarKem to achieve a better trade-off between decryption failure probability and security than a fixed BW construction.

\subsection{Main Contributions}\label{sec:contributions}

The main contributions of this document are:

\begin{enumerate}[leftmargin=2em]
\item \textbf{A new \lipPrb-based KEM framework} using polar-code-defined lattices, with a complete formal specification including key generation, encapsulation, and decapsulation algorithms.
\item \textbf{Three concrete parameter sets} (\textsf{PolarKEM-128}, \textsf{PolarKEM-256}, \textsf{PolarKEM-512}) targeting NIST security levels 1, 3, and 5, with detailed security analysis and attack estimates.
\item \textbf{Security proofs} (in the random oracle model) establishing \INDCPA security for the core KEM and \INDCCA security after the Fujisaki-Okamoto transform.
\item \textbf{Comprehensive performance analysis} including key sizes, ciphertext sizes, computational complexity, and comparison with other post-quantum KEMs.
\end{enumerate}

\subsection{Document Organization}\label{sec:organization}

The remainder of this document is organized as follows. Section~\ref{sec:prelim} reviews mathematical preliminaries. Section~\ref{sec:design} provides a high-level design overview. Section~\ref{sec:specification} gives the formal specification of \polarKem. Section~\ref{sec:security} analyzes security. Section~\ref{sec:performance} evaluates performance. Section~\ref{sec:implementation} describes implementation notes. Sections~\ref{sec:features} and~\ref{sec:limitations} state features, limitations, and open problems. References follow, with appendices providing mathematical background, proof details, alternative designs, and parameter selection methodology.

%%%%%%%%%%%%%%%%%%%%%%%%%%%%%%%%%%%%%%%%%%%%%%%%%%%%%%%%%%%%%%%%%%%%%%%%%%%%%%%%
% PRELIMINARIES
%%%%%%%%%%%%%%%%%%%%%%%%%%%%%%%%%%%%%%%%%%%%%%%%%%%%%%%%%%%%%%%%%%%%%%%%%%%%%%%%
\section{Preliminaries}\label{sec:prelim}

\subsection{Notation}\label{sec:notation}

We use the following notation throughout this document:

\begin{itemize}[leftmargin=2em]
\item $\bit = \{0, 1\}$. For $n \in \mathbb{N}$, $\bit^n$ is the set of $n$-bit strings.
\item $\Z_q = \Z / q\Z$ denotes the ring of integers modulo $q$.
\item Vectors are denoted by bold lowercase letters $\mathbf{v}$; matrices by bold uppercase letters $\mathbf{M}$.
\item $\mathbf{v}^T$ denotes the transpose of $\mathbf{v}$. The inner product is $\langle \mathbf{u}, \mathbf{v} \rangle = \mathbf{u}^T \mathbf{v}$.
\item $\norm{\mathbf{v}}$ denotes the Euclidean ($\ell_2$) norm of $\mathbf{v}$.
\item $\GL_n(\Z)$ is the group of $n \times n$ unimodular integer matrices.
\item $\mathcal{O}_n(\R)$ is the group of $n \times n$ real orthogonal matrices.
\item $\lambda$ denotes the security parameter.
\item $\ poly(\lambda)$ denotes a polynomial in $\lambda$.
\item $\negl(\lambda)$ denotes a negligible function in $\lambda$.
\item $x \sample S$ denotes sampling $x$ uniformly at random from a finite set $S$.
\item $x \gets \mathcal{D}$ denotes sampling from a distribution $\mathcal{D}$.
\item $\lfloor x \rceil$ denotes rounding $x$ to the nearest integer.
\end{itemize}

\subsection{Lattices and Quadratic Forms}\label{sec:lattices}

\begin{definition}[Lattice]
A \emph{lattice} $\Lambda \subset \R^n$ is a discrete additive subgroup of $\R^n$. Equivalently, $\Lambda$ is the set of all integer linear combinations of a set of linearly independent basis vectors $\{\mathbf{b}_1, \ldots, \mathbf{b}_n\}$:
\[
\Lambda = \mathcal{L}(\mathbf{B}) = \left\{ \sum_{i=1}^{n} z_i \mathbf{b}_i : z_i \in \Z \right\},
\]
where $\mathbf{B} = [\mathbf{b}_1 | \cdots | \mathbf{b}_n] \in \R^{n \times n}$ is the basis matrix.
\end{definition}

\begin{definition}[Gram matrix]
The \emph{Gram matrix} of a lattice $\Lambda$ with basis $\mathbf{B}$ is $\mathbf{G} = \mathbf{B}^T \mathbf{B}$. The \emph{Gram--Schmidt orthogonalization} $\widetilde{\mathbf{B}} = [\widetilde{\mathbf{b}}_1 | \cdots | \widetilde{\mathbf{b}}_n]$ is defined recursively by
\[
\widetilde{\mathbf{b}}_i = \mathbf{b}_i - \sum_{j=1}^{i-1} \mu_{i,j} \widetilde{\mathbf{b}}_j, \qquad \mu_{i,j} = \frac{\langle \mathbf{b}_i, \widetilde{\mathbf{b}}_j \rangle}{\langle \widetilde{\mathbf{b}}_j, \widetilde{\mathbf{b}}_j \rangle}.
\]
\end{definition}

\begin{definition}[Quadratic form]
A \emph{positive definite quadratic form} $Q: \R^n \to \R$ is a function $Q(\mathbf{x}) = \mathbf{x}^T \mathbf{G} \mathbf{x}$ where $\mathbf{G}$ is a symmetric positive definite matrix. The \emph{minimum} of $Q$ is $\min(Q) = \min_{\mathbf{x} \in \Z^n \setminus \{\mathbf{0}\}} Q(\mathbf{x})$.
\end{definition}

\begin{definition}[Determinant]
The \emph{determinant} of a lattice $\Lambda = \mathcal{L}(\mathbf{B})$ is $\det(\Lambda) = \sqrt{\det(\mathbf{B}^T \mathbf{B})}$.
\end{definition}

\subsection{The Lattice Isomorphism Problem}\label{sec:lip}

\begin{definition}[Lattice Isomorphism]
Two lattices $\Lambda_1, \Lambda_2 \subset \R^n$ are \emph{isomorphic} (written $\Lambda_1 \cong \Lambda_2$) if there exists an orthogonal matrix $\mathbf{O} \in \mathcal{O}_n(\R)$ such that $\Lambda_2 = \mathbf{O} \cdot \Lambda_1$. Equivalently, $\Lambda_1 \cong \Lambda_2$ if and only if they share the same Gram matrix up to unimodular equivalence.
\end{definition}

\begin{definition}[Lattice Isomorphism Problem (\lipPrb)]\label{def:lip}
Given bases $\mathbf{B}_1, \mathbf{B}_2$ of two isomorphic lattices $\Lambda_1 \cong \Lambda_2$, find an orthogonal matrix $\mathbf{O} \in \mathcal{O}_n(\R)$ and a unimodular matrix $\mathbf{U} \in \GL_n(\Z)$ such that $\mathbf{B}_2 = \mathbf{O} \cdot \mathbf{B}_1 \cdot \mathbf{U}$.
\end{definition}

\begin{assumption}[Hardness of \lipPrb]\label{ass:lip}
For certain families of lattices (including polar-code-defined lattices), there is no probabilistic polynomial-time algorithm that solves the Lattice Isomorphism Problem with non-negligible probability.
\end{assumption}

\begin{assumption}[\polarLIP Assumption]\label{ass:polar-lip}
Let $\mathcal{P}_n$ be the family of polar lattices of dimension $n$. The \emph{\polarLIP} assumption states that for $\Lambda \sample \mathcal{P}_n$ and a random orthogonal $\mathbf{O} \sample \mathcal{O}_n(\R)$, given the basis of $\Lambda' = \mathbf{O} \cdot \Lambda$, it is hard to recover $\mathbf{O}$ (or any isometry).
\end{assumption}

\subsection{Polar Codes and Polar Lattices}\label{sec:polar}

\subsubsection{Polar Codes}

Polar codes, introduced by Ar{\i}kan~\cite{Arikan2009}, are a family of linear block codes that achieve the symmetric capacity of binary-input discrete memoryless channels under successive cancellation decoding. A polar code of length $N = 2^m$ and dimension $K$ is defined by a \emph{frozen set} $\mathcal{F} \subset \{0, 1, \ldots, N-1\}$ of size $N - K$.

The polar transform is defined by the $N \times N$ matrix $\mathbf{G}_N = \mathbf{B}_N \mathbf{F}^{\otimes m}$, where $\mathbf{F} = \begin{psmallmatrix} 1 & 0 \\ 1 & 1 \end{psmallmatrix}$, $\otimes m$ denotes the $m$-th Kronecker power, and $\mathbf{B}_N$ is the bit-reversal permutation matrix.

For a message $\mathbf{u} \in \bit^N$ where $u_i = 0$ for all $i \in \mathcal{F}$, the codeword is $\mathbf{c} = \mathbf{u} \cdot \mathbf{G}_N$.

\subsubsection{Polar Lattices}

\begin{definition}[Polar lattice via Construction D]\label{def:polar-lattice}
Let $N = 2^m$ and let $\mathcal{F}_0 \subseteq \mathcal{F}_1 \subseteq \cdots \subseteq \mathcal{F}_{L-1} \subseteq \{0, \ldots, N-1\}$ be a chain of frozen sets. For $\ell = 0, \ldots, L-1$, let $C_\ell$ be the polar code with frozen set $\mathcal{F}_\ell$. The \emph{polar lattice} $\Lambda_{\text{polar}}$ is defined by
\[
\Lambda_{\text{polar}} = C_0 + 2C_1 + 4C_2 + \cdots + 2^{L-1}C_{L-1} + 2^L \Z^N.
\]
This is an instance of \emph{Construction D}.
\end{definition}

The volume of $\Lambda_{\text{polar}}$ is $\det(\Lambda_{\text{polar}}) = 2^{N \cdot L - \sum_{\ell=0}^{L-1} K_\ell}$, where $K_\ell = N - |\mathcal{F}_\ell|$ is the dimension of $C_\ell$.

\subsubsection{Nested Structure and Rate Progression}

In \polarKem, we use a rate progression where each level $\ell$ has rate $R_\ell = 2^{-\ell}$. This gives:
\[
K_\ell = N \cdot R_\ell = \frac{N}{2^\ell}.
\]

The nested structure $\mathcal{F}_0 \subseteq \mathcal{F}_1 \subseteq \cdots \subseteq \mathcal{F}_{L-1}$ ensures that the codes are compatible with Construction D.

\subsubsection{Minimum Distance}

The minimum distance of the polar lattice is determined by the minimum distance of the base code $C_0$ and the construction level $L$.

\begin{lemma}[Minimum Distance of Polar Lattice]\label{lem:min-dist}
Let $\Lambda$ be a polar lattice via Construction~D with $L$ levels. The squared minimum distance satisfies:
\[
d(\Lambda)^2 \geq \min\left(d(C_0), 4 \cdot d(C_1), \ldots, 4^{L-1} \cdot d(C_{L-1}), 4^L\right),
\]
where $d(C_\ell)$ is the minimum Hamming distance of code $C_\ell$.
\end{lemma}

\subsection{Construction D and Construction D$'$}\label{sec:construction-d}

\begin{definition}[Construction D]\label{def:construction-d}
Let $C_0 \supseteq C_1 \supseteq \cdots \supseteq C_{L-1}$ be a chain of binary linear codes of length $N$, where each $C_\ell$ has dimension $K_\ell$. Let $\{\mathbf{g}_1, \ldots, \mathbf{g}_N\}$ be vectors in $\R^N$ such that $\mathbf{g}_1, \ldots, \mathbf{g}_{K_\ell}$ reduced modulo 2 generate $C_\ell$. Then Construction~D produces the lattice:
\[
\Lambda = \left\{ \sum_{\ell=0}^{L-1} \sum_{j=1}^{K_\ell} 2^\ell u_{\ell,j} \mathbf{g}_j + 2^L \mathbf{z} : u_{\ell,j} \in \{0,1\}, \mathbf{z} \in \Z^N \right\}.
\]
\end{definition}

\begin{definition}[Construction D$'$]\label{def:construction-d-prime}
Construction D$'$ is a variant where the codes are nested in the opposite direction ($C_0 \subseteq C_1 \subseteq \cdots \subseteq C_{L-1}$) and the lattice is constructed using a different scaling pattern. Construction D$'$ is not used in \polarKem but is mentioned for comparison with Barnes-Wall lattice constructions.
\end{definition}

Polar lattices in \polarKem use Construction~D with the chain $C_0 \supseteq C_1 \supseteq \cdots \supseteq C_{L-1}$ obtained from polar codes with nested frozen sets.

\subsection{Bounded-Distance Decoding, CVP, and uSVP}\label{sec:bdd}

\begin{definition}[Bounded-Distance Decoding (\BDD)]\label{def:bdd}
Given a lattice $\Lambda \subset \R^n$ and a target $\mathbf{t} = \mathbf{v} + \mathbf{e}$ where $\mathbf{v} \in \Lambda$ and $\norm{\mathbf{e}} < \rho$ for some decoding radius $\rho$, the \BDD problem is to recover $\mathbf{v}$.
\end{definition}

\begin{definition}[Closest Vector Problem (\CVP)]\label{def:cvp}
Given a lattice $\Lambda$ with basis $\mathbf{B}$ and a target $\mathbf{t} \in \R^n$, find the lattice vector $\mathbf{v} \in \Lambda$ minimizing $\norm{\mathbf{t} - \mathbf{v}}$.
\end{definition}

\begin{definition}[Unique Shortest Vector Problem (\usvp)]\label{def:usvp}
Given a lattice $\Lambda$ where $\lambda_1(\Lambda) \leq \gamma \cdot \lambda_2(\Lambda)$ for some gap $\gamma > 1$, find a shortest non-zero vector of $\Lambda$.
\end{definition}

\begin{lemma}[BDD reduces to \usvp for small error]\label{lem:bdd-usvp}
If $\norm{\mathbf{e}} < \lambda_1(\Lambda) / 2$, then \BDD reduces to \CVP, and for suitable approximation factors, \CVP reduces to \usvp via the Kannan embedding.
\end{lemma}

\subsection{Randomness Extractors}\label{sec:extractors}

\begin{definition}[Strong Extractor]\label{def:extractor}
A function $\Ext: \bit^n \times \bit^d \to \bit^m$ is a \emph{$(k, \epsilon)$-strong extractor} if for all random variables $X$ over $\bit^n$ with min-entropy $H_\infty(X) \geq k$, the distribution $(S, \Ext(X, S))$ is $\epsilon$-close to uniform, where $S \sample \bit^d$ is the seed.
\end{definition}

In \polarKem, we use a hash-based extractor modeled as a random oracle. The Leftover Hash Lemma ensures that universal hash functions yield good extractors.

\subsection{KEM Security Notions}\label{sec:security-notions}

A Key Encapsulation Mechanism (KEM) is a triple $\mathsf{KEM} = (\mathsf{KeyGen}, \mathsf{Encaps}, \mathsf{Decaps})$:

\begin{itemize}[leftmargin=2em]
\item $(\pk, \sk) \gets \mathsf{KeyGen}(1^\lambda)$: generates public key $\pk$ and secret key $\sk$.
\item $(\mathbf{c}, K) \gets \mathsf{Encaps}(\pk)$: generates ciphertext $\mathbf{c}$ and shared secret $K$.
\item $K \gets \mathsf{Decaps}(\sk, \mathbf{c})$: recovers shared secret $K$ or returns $\bot$.
\end{itemize}

\subsubsection{IND-CPA Security}

\begin{definition}[IND-CPA]\label{def:indcpa}
A KEM is \emph{IND-CPA secure} if for all probabilistic polynomial-time adversaries $\mathcal{A}$,
\[
\left|\Pr\left[\mathcal{A}(\pk, \mathbf{c}, K_0) = 1\right] - \Pr\left[\mathcal{A}(\pk, \mathbf{c}, K_1) = 1\right]\right| \leq \negl(\lambda),
\]
where $(\pk, \sk) \gets \mathsf{KeyGen}(1^\lambda)$, $(\mathbf{c}, K_0) \gets \mathsf{Encaps}(\pk)$, $K_1 \sample \bit^{|K_0|}$, and $b \sample \{0, 1\}$.
\end{definition}

\subsubsection{IND-CCA2 Security}

\begin{definition}[IND-CCA2]\label{def:indcca}
A KEM is \emph{IND-CCA2 secure} if for all probabilistic polynomial-time adversaries $\mathcal{A}$ with access to a decapsulation oracle (that cannot be called on the challenge ciphertext),
\[
\left|\Pr[\text{Exp}_{\mathcal{A}}^{\text{IND-CCA2}}(1^\lambda) = 1] - \frac{1}{2}\right| \leq \negl(\lambda).
\]
\end{definition}


%%%%%%%%%%%%%%%%%%%%%%%%%%%%%%%%%%%%%%%%%%%%%%%%%%%%%%%%%%%%%%%%%%%%%%%%%%%%%%%%
% DESIGN OVERVIEW
%%%%%%%%%%%%%%%%%%%%%%%%%%%%%%%%%%%%%%%%%%%%%%%%%%%%%%%%%%%%%%%%%%%%%%%%%%%%%%%%
\section{Design Overview}\label{sec:design}

\subsection{High-Level Construction}\label{sec:high-level}

\polarKem follows the blueprint of LIP-based encryption. The public key is a disguised (rotated) version of a polar lattice basis; the secret key is the orthogonal transformation that ``undisguises'' it. Encryption works as follows:

\begin{enumerate}[leftmargin=2em]
\item The encapsulator samples a short error vector $\mathbf{e}$ with bounded norm.
\item The encapsulator chooses a random message $\mathbf{m}$ and encodes it to a lattice vector $\mathbf{v}$ using the \emph{public} (rotated) lattice.
\item The ciphertext is $\mathbf{c} = \mathbf{v} + \mathbf{e}$.
\item The shared secret is derived from $\mathbf{m}$ via an extractor.
\end{enumerate}

Decapsulation uses the secret orthogonal transformation $\mathbf{O}$ to map $\mathbf{c}$ to the non-rotated lattice:
\[
\mathbf{O}^{-1} \mathbf{c} = \mathbf{O}^{-1} \mathbf{v} + \mathbf{O}^{-1} \mathbf{e} = \mathbf{v}' + \mathbf{e}',
\]
where $\mathbf{v}'$ is a vector in the original polar lattice. Because the polar lattice admits efficient bounded-distance decoding, the decapsulator recovers $\mathbf{v}'$ (and thus $\mathbf{m}$) as long as $\norm{\mathbf{e}'} = \norm{\mathbf{e}}$ is within the decoding radius.

\subsection{Core KEM (IND-CPA)}\label{sec:core-kem}

The \INDCPA-secure \CoreKEM consists of three algorithms:

\begin{itemize}[leftmargin=2em]
\item $\CoreKEM.\mathsf{KeyGen}(1^\lambda)$: Generates a polar lattice basis $\mathbf{B}$, samples a random orthogonal matrix $\mathbf{O}$, computes the public basis $\mathbf{B}_{\pk} = \mathbf{O} \cdot \mathbf{B}$, and outputs $(\pk = \mathbf{B}_{\pk}, \sk = \mathbf{O})$.
\item $\CoreKEM.\mathsf{Encaps}(\pk)$: Samples message $\mathbf{m}$, encodes it to $\mathbf{v} \in \mathcal{L}(\mathbf{B}_{\pk})$, samples short error $\mathbf{e}$, computes $\mathbf{c} = \mathbf{v} + \mathbf{e}$, and outputs $(\mathbf{c}, K = \Extract(\mathbf{m}))$.
\item $\CoreKEM.\mathsf{Decaps}(\sk, \mathbf{c})$: Computes $\mathbf{c}' = \mathbf{O}^{-1} \mathbf{c}$, decodes $\mathbf{c}'$ to recover $\mathbf{m}'$ using the polar lattice SC decoder, and outputs $K = \Extract(\mathbf{m}')$.
\end{itemize}

\subsection{CCA2 Transformation}\label{sec:cca-transform}

\polarKem applies the Fujisaki-Okamoto (\FO) transform with implicit rejection to achieve \INDCCA security in the \QROM. The \FO transform:

\begin{enumerate}[leftmargin=2em]
\item Derives both the message and random coins from a single seed using a pseudorandom generator.
\item Re-encrypts during decapsulation to verify ciphertext integrity.
\item Uses implicit rejection: on verification failure, derives a pseudorandom shared secret from the secret key and ciphertext instead of returning $\bot$.
\end{enumerate}

\subsection{Design Rationale}\label{sec:design-rationale}

\subsubsection{Choice of Polar Lattices}

Polar lattices were chosen for several reasons:
\begin{enumerate}[leftmargin=2em]
\item \textbf{Efficient decoding:} SC decoding runs in $\mathcal{O}(N \log N)$ time.
\item \textbf{Controllable DFP:} The decoding radius is precisely determined by the polar code parameters, allowing the decryption failure probability to be made cryptographically negligible.
\item \textbf{Structural diversity:} Polar lattices offer a different structural profile than Barnes-Wall or orthogonal lattices, diversifying the LIP assumption landscape.
\end{enumerate}

\subsubsection{Choice of Modulus}

The modulus $q$ is chosen to enable efficient arithmetic while ensuring that the ring structure does not introduce algebraic attacks. We use $q = 12289$ for \textsf{PolarKEM-128} and \textsf{PolarKEM-256} (matching the Kyber/ML-KEM modulus for hardware compatibility) and $q = 18433$ for \textsf{PolarKEM-512} to provide additional security margin at the highest security level.

\subsubsection{Error Distribution}

The error distribution is a narrow discrete Gaussian with small standard deviation. The error bound $B = 1$ (i.e., each error component is in $\{-1, 0, 1\}$) ensures that errors are short enough for reliable decoding while maximizing the entropy of the error.

%%%%%%%%%%%%%%%%%%%%%%%%%%%%%%%%%%%%%%%%%%%%%%%%%%%%%%%%%%%%%%%%%%%%%%%%%%%%%%%%
% FORMAL SPECIFICATION
%%%%%%%%%%%%%%%%%%%%%%%%%%%%%%%%%%%%%%%%%%%%%%%%%%%%%%%%%%%%%%%%%%%%%%%%%%%%%%%%
\section{Formal Specification of \polarKem}\label{sec:specification}

\subsection{Parameters}\label{sec:parameters}

Table~\ref{tab:params} defines the parameter sets for \polarKem.

\begin{table}[H]
\centering
\caption{Parameter sets for \polarKem.}
\label{tab:params}
\begin{tabular}{@{}lccc@{}}
\toprule
\textbf{Parameter} & \textbf{PolarKEM-128} & \textbf{PolarKEM-256} & \textbf{PolarKEM-512} \\
\midrule
$N$ (dimension) & 512 & 1024 & 2048 \\
$q$ (modulus) & 12289 & 12289 & 18433 \\
$L$ (levels) & 5 & 5 & 6 \\
Error bound $B$ & 1 & 1 & 1 \\
Decoding radius & 15.0 & 31.0 & 63.0 \\
Shared secret length & 128 bits & 256 bits & 512 bits \\
$\delta$ (decryption failure prob.) & $< 2^{-128}$ & $< 2^{-256}$ & $< 2^{-512}$ \\
Classical security & $\sim$150 bits & $\sim$296 bits & $\sim$618 bits \\
Quantum security & $\sim$134 bits & $\sim$269 bits & $\sim$561 bits \\
\bottomrule
\end{tabular}
\end{table}

\subsubsection{Parameter Constraints}

The parameters must satisfy the following constraints:
\begin{enumerate}[leftmargin=2em]
\item The dimension $N = 2^m$ is a power of two for the polar transform.
\item The modulus $q$ is a prime satisfying $q \equiv 1 \pmod{2N}$ to support NTT-based polynomial multiplication.
\item The number of levels $L$ satisfies $2^L \leq q$ to ensure correct Construction~D encoding.
\item The decoding radius $\rho$ satisfies $\rho < \lambda_1(\Lambda) / 2$ where $\Lambda$ is the polar lattice.
\end{enumerate}

\subsection{Core KEM: Key Generation}\label{sec:core-keygen}

Algorithm~\ref{alg:core-keygen} specifies the \INDCPA key generation algorithm.

\begin{algorithm}[H]
\caption{$\CoreKEM.\mathsf{KeyGen}()$}
\label{alg:core-keygen}
\begin{algorithmic}[1]
\REQUIRE Parameters $(N, q, L)$
\ENSURE Public key $\pk$, secret key $\sk$
\STATE $m \gets \log_2 N$
\STATE Construct polar code chain $C_0 \supseteq C_1 \supseteq \cdots \supseteq C_{L-1}$ with frozen sets $\mathcal{F}_0 \subseteq \mathcal{F}_1 \subseteq \cdots \subseteq \mathcal{F}_{L-1}$
\STATE Build polar lattice basis $\mathbf{B} \in \Z^{N \times N}$ via Construction~D from the code chain
\STATE $\mathbf{B}_{\text{red}} \gets \textsc{LLL}(\mathbf{B})$ \COMMENT{Reduce for good decoding properties}
\STATE Sample orthogonal matrix $\mathbf{O} \sample \mathcal{O}_N(\R)$ uniformly at random
\STATE $\mathbf{B}_{\pk} \gets \mathbf{O} \cdot \mathbf{B}_{\text{red}}$
\STATE $\pk \gets \mathbf{B}_{\pk}$
\STATE $\sk \gets \mathbf{O}$
\RETURN $(\pk, \sk)$
\end{algorithmic}
\end{algorithm}

\subsection{Core KEM: Encapsulation}\label{sec:core-encaps}

Algorithm~\ref{alg:core-encaps} specifies the \INDCPA encapsulation algorithm.

\begin{algorithm}[H]
\caption{$\CoreKEM.\mathsf{Encaps}(\pk)$}
\label{alg:core-encaps}
\begin{algorithmic}[1]
\REQUIRE Public key $\pk = \mathbf{B}_{\pk}$
\ENSURE Ciphertext $\mathbf{c}$, shared secret $K$
\STATE $\mathbf{m} \sample \bit^K$ \COMMENT{Sample random message}
\STATE $\mathbf{v} \gets \mathsf{Embed}(\mathbf{m}, \mathbf{B}_{\pk})$ \COMMENT{Encode message to lattice vector}
\STATE $\mathbf{e} \sample \mathcal{D}_{\text{err}}$ \COMMENT{Sample short error vector}
\STATE $\mathbf{c} \gets \mathbf{v} + \mathbf{e} \pmod{q}$
\STATE $K \gets \Extract(\mathbf{m})$
\RETURN $(\mathbf{c}, K)$
\end{algorithmic}
\end{algorithm}

The $\mathsf{Embed}$ function maps a message $\mathbf{m} \in \bit^K$ to a lattice vector $\mathbf{v} \in \mathcal{L}(\mathbf{B}_{\pk})$ by interpreting $\mathbf{m}$ as the information bits of the polar code chain and applying the polar transform.

\subsection{Core KEM: Decapsulation}\label{sec:core-decaps}

Algorithm~\ref{alg:core-decaps} specifies the \INDCPA decapsulation algorithm.

\begin{algorithm}[H]
\caption{$\CoreKEM.\mathsf{Decaps}(\sk, \mathbf{c})$}
\label{alg:core-decaps}
\begin{algorithmic}[1]
\REQUIRE Secret key $\sk = \mathbf{O}$, ciphertext $\mathbf{c}$
\ENSURE Shared secret $K$
\STATE $\mathbf{c}' \gets \mathbf{O}^{-1} \cdot \mathbf{c} \pmod{q}$ \COMMENT{Rotate to secret lattice}
\STATE $\mathbf{c}'' \gets \mathbf{c}' \bmod 2^L$ \COMMENT{Reduce to fundamental region}
\STATE $\widehat{\mathbf{m}} \gets \polarScDec(\mathbf{c}'', \{\mathcal{F}_\ell\}_{\ell=0}^{L-1})$ \COMMENT{SC decode}
\STATE $K \gets \Extract(\widehat{\mathbf{m}})$
\RETURN $K$
\end{algorithmic}
\end{algorithm}

\subsection{Polar Lattice Construction Details}\label{sec:polar-construction}

\subsubsection{Frozen Set Selection}

The frozen sets are selected using the polar code design rule for the Binary Memoryless Symmetric (BMS) channel with erasure probability corresponding to the error distribution. For each level $\ell$, we compute the Bhattacharyya parameters $\{Z_i^{(\ell)}\}_{i=0}^{N-1}$ and freeze the $N - K_\ell$ indices with the highest values.

The rate progression for \polarKem is:
\[
R_\ell = 2^{-\ell} \quad \Rightarrow \quad K_\ell = \frac{N}{2^\ell}.
\]

\subsubsection{Construction D Implementation}

The polar lattice basis $\mathbf{B}$ is constructed as follows:

\begin{enumerate}[leftmargin=2em]
\item For each level $\ell \in \{0, \ldots, L-1\}$:
  \begin{enumerate}
  \item Generate the polar code generator matrix $\mathbf{G}_\ell$ for code $C_\ell$.
  \item Extract rows corresponding to information bits (non-frozen indices).
  \end{enumerate}
\item Stack the scaled generator matrices to form $\mathbf{B}_{\text{code}} = [\mathbf{G}_0^T | 2\mathbf{G}_1^T | \cdots | 2^{L-1}\mathbf{G}_{L-1}^T]^T$.
\item Append $2^L \mathbf{I}_N$ to handle the integer lattice component.
\item Apply LLL reduction to obtain a basis suitable for decoding.
\end{enumerate}

The formal basis matrix is:
\[
\mathbf{B} = \begin{bmatrix}
\mathbf{G}_0^T & 2\mathbf{G}_1^T & \cdots & 2^{L-1}\mathbf{G}_{L-1}^T & 2^L \mathbf{I}_N
\end{bmatrix}^T
\]
where each $\mathbf{G}_\ell$ contains only the rows of the polar transform corresponding to information (non-frozen) positions at level $\ell$.

\subsubsection{Minimum Distance Formula}

For the rate progression $R_\ell = 2^{-\ell}$, the minimum distance of the polar lattice satisfies:
\[
d(\Lambda)^2 \geq 4^L
\]
since the coarsest code $C_{L-1}$ has rate $R_{L-1} = 2^{-(L-1)}$ and the minimum distance of the full construction is bounded below by $4^L$ (from the $\Z^N$ component at level $L$).

The actual minimum distance may be larger depending on the minimum distances of the individual polar codes. For the parameter sets in Table~\ref{tab:params}, the minimum distances are:

\begin{table}[H]
\centering
\caption{Minimum distance estimates for \polarKem parameter sets.}
\label{tab:min-dist}
\begin{tabular}{@{}lccc@{}}
\toprule
 & \textbf{PolarKEM-128} & \textbf{PolarKEM-256} & \textbf{PolarKEM-512} \\
\midrule
$d^2(\Lambda)$ (lower bound) & $4^5 = 1024$ & $4^5 = 1024$ & $4^6 = 4096$ \\
Estimated $d^2(\Lambda)$ & $\approx 1536$ & $\approx 3072$ & $\approx 8192$ \\
Decoding radius $\rho$ & 15.0 & 31.0 & 63.0 \\
$\rho / d(\Lambda)$ & $\approx 0.383$ & $\approx 0.559$ & $\approx 0.696$ \\
\bottomrule
\end{tabular}
\end{table}

\subsection{Successive Cancellation Decoder Specification}\label{sec:sc-decoder}

Algorithm~\ref{alg:sc-decode} specifies the successive cancellation decoder for polar lattices.

\begin{algorithm}[H]
\caption{$\polarScDec(\mathbf{y}, \{\mathcal{F}_\ell\}_{\ell=0}^{L-1})$}
\label{alg:sc-decode}
\begin{algorithmic}[1]
\REQUIRE Received vector $\mathbf{y} \in \R^N$, frozen sets $\{\mathcal{F}_\ell\}$
\ENSURE Decoded message $\widehat{\mathbf{m}}$
\STATE $\widehat{\mathbf{u}} \gets \mathbf{0} \in \Z^N$
\STATE $m \gets \log_2 N$
\FOR{$i = 0$ \textbf{to} $N - 1$}
    \STATE $\ell^* \gets \min\{\ell : i \in \mathcal{F}_\ell\}$ \COMMENT{Find coarsest level containing $i$}
    \IF{$\ell^*$ exists}
        \STATE $\widehat{u}_i \gets 0$ \COMMENT{Frozen bit}
    \ELSE
        \STATE $L_i \gets L - 1$ \COMMENT{Information bit at finest level}
        \STATE Compute log-likelihood ratio: $\text{LLR}_i \gets \log \frac{p(y_i | u_i = 0)}{p(y_i | u_i = 1)}$
        \STATE $\widehat{u}_i \gets \CTSelect(0, 1, \text{LLR}_i \geq 0)$
    \ENDIF
\ENDFOR
\STATE \COMMENT{Propagate through polar transform}
\STATE $\widehat{\mathbf{m}} \gets \widehat{\mathbf{u}} \cdot \mathbf{G}_N \pmod{2^L}$
\RETURN $\widehat{\mathbf{m}}$
\end{algorithmic}
\end{algorithm}

The SC decoder processes each bit sequentially. For each position $i$, the decoder computes a log-likelihood ratio (LLR) based on the received values and previously decoded bits. If $i$ is a frozen index at level $\ell$, the bit is set to 0; otherwise, it is decoded using the LLR.

\subsubsection{Multi-Level Decoding}

For Construction~D with $L$ levels, the SC decoder operates level by level:

\begin{enumerate}[leftmargin=2em]
\item Start with the finest level ($\ell = L-1$), decoding modulo 2.
\item Subtract the decoded contribution and divide by 2.
\item Proceed to the next level, decoding modulo 2 again.
\item Continue until all levels are processed.
\end{enumerate}

This multi-level approach ensures that each level's code is decoded in its appropriate domain.

\subsection{Error Distribution}\label{sec:error-dist}

The error distribution $\mathcal{D}_{\text{err}}$ is a discrete distribution over $\Z^N$ where each component $e_i$ is sampled independently from a narrow distribution centered at 0.

\begin{definition}[Error Distribution]\label{def:error-dist}
For \polarKem, the error distribution is:
\[
e_i \sample \begin{cases}
0 & \text{with probability } p_0, \\
+1 & \text{with probability } p_1, \\
-1 & \text{with probability } p_{-1},
\end{cases}
\]
where $p_0 + p_1 + p_{-1} = 1$ and $p_1 = p_{-1} = (1 - p_0)/2$. The default parameter $p_0 = 1/2$ gives a centered binomial-like distribution with variance $1/2$.
\end{definition}

The expected squared norm of the error is:
\[
\mathbb{E}[\norm{\mathbf{e}}^2] = N \cdot \mathbb{E}[e_i^2] = N \cdot (1 - p_0).
\]

For $p_0 = 1/2$, we have $\mathbb{E}[\norm{\mathbf{e}}^2] = N/2$.

The probability that $\norm{\mathbf{e}} \leq \rho$ (the decoding success probability) is computed via the Gaussian approximation to the binomial distribution and is designed to satisfy the \DFP bounds in Table~\ref{tab:params}.

\subsection{Extractor and Key Derivation Function}\label{sec:extractor-kdf}

\subsubsection{Randomness Extractor}

The extractor $\Extract: \bit^N \to \bit^\ell$ is defined as:
\[
\Extract(\mathbf{m}) = \textsc{SHA3-256}(\mathbf{m} \| \text{``PolarKEM-v1''}),
\]
where $\ell$ is the shared secret length (128, 256, or 512 bits). For 512-bit shared secrets, the output is extended using $\textsc{SHAKE-256}$.

\subsubsection{Key Derivation Function}

The KDF used in the \FO transform is:
\[
\KDF(s) = \textsc{SHAKE-256}(s \| \text{``PolarKEM-FO''}).
\]

\subsection{CCA Transform: Fujisaki-Okamoto with Implicit Rejection}\label{sec:cca-fo}

Algorithm~\ref{alg:kem-keygen}, Algorithm~\ref{alg:kem-encaps}, and Algorithm~\ref{alg:kem-decaps} specify the full \INDCCA-secure \polarKem.

\begin{algorithm}[H]
\caption{$\polarKem.\mathsf{KeyGen}(1^\lambda)$}
\label{alg:kem-keygen}
\begin{algorithmic}[1]
\REQUIRE Security parameter $\lambda$
\ENSURE Public key $\pk$, secret key $\sk$
\STATE $(\pk_0, \sk_0) \gets \CoreKEM.\mathsf{KeyGen}(1^\lambda)$
\STATE $s \sample \bit^{\lambda}$ \COMMENT{Implicit rejection seed}
\STATE $\pk \gets \pk_0$
\STATE $\sk \gets (\sk_0, s, \pk)$
\RETURN $(\pk, \sk)$
\end{algorithmic}
\end{algorithm}

\begin{algorithm}[H]
\caption{$\polarKem.\mathsf{Encaps}(\pk)$}
\label{alg:kem-encaps}
\begin{algorithmic}[1]
\REQUIRE Public key $\pk$
\ENSURE Ciphertext $\mathbf{c}$, shared secret $K$
\STATE $\mu \sample \bit^{\lambda}$ \COMMENT{Random seed}
\STATE $(\mathbf{m}, \mathbf{r}) \gets \KDF(\mu)$ \COMMENT{Derive message and random coins}
\STATE $\mathbf{v} \gets \mathsf{Embed}(\mathbf{m}, \pk)$
\STATE $\mathbf{e} \gets \mathsf{SampleErr}(\mathbf{r})$
\STATE $\mathbf{c} \gets \mathbf{v} + \mathbf{e} \pmod{q}$
\STATE $K \gets \Extract(\mathbf{m})$
\RETURN $(\mathbf{c}, K)$
\end{algorithmic}
\end{algorithm}

\begin{algorithm}[H]
\caption{$\polarKem.\mathsf{Decaps}(\sk, \mathbf{c})$}
\label{alg:kem-decaps}
\begin{algorithmic}[1]
\REQUIRE Secret key $\sk = (\sk_0, s, \pk)$, ciphertext $\mathbf{c}$
\ENSURE Shared secret $K$
\STATE $\widehat{\mathbf{m}} \gets \CoreKEM.\mathsf{Decaps}(\sk_0, \mathbf{c})$ \COMMENT{Core decapsulation}
\STATE $\widehat{\mathbf{v}} \gets \mathsf{Embed}(\widehat{\mathbf{m}}, \pk)$
\STATE $\widehat{\mathbf{e}} \gets \mathbf{c} - \widehat{\mathbf{v}} \pmod{q}$
\STATE \textbf{if} $\norm{\widehat{\mathbf{e}}} > \rho$ \textbf{then} \COMMENT{Check decoding radius}
    \STATE \quad $K \gets \KDF(s \| \mathbf{c})$ \COMMENT{Implicit rejection}
\STATE \textbf{else}
    \STATE \quad $K \gets \Extract(\widehat{\mathbf{m}})$
\STATE \textbf{end if}
\RETURN $K$
\end{algorithmic}
\end{algorithm}

\textbf{Implicit Rejection:} If the ciphertext fails the validity check (line~4), the decapsulation algorithm does not return $\bot$. Instead, it derives a pseudorandom shared secret from the secret seed $s$ and the ciphertext $\mathbf{c}$. This ensures that timing attacks cannot distinguish valid from invalid ciphertexts.

\subsection{Constant-Time Requirements}\label{sec:constant-time}

All operations in \polarKem must be implemented in constant time with respect to secret data. This includes:

\begin{enumerate}[leftmargin=2em]
\item The SC decoder must not branch on secret bit values.
\item The orthogonal transformation must not leak timing information.
\item The decapsulation must not branch on whether the ciphertext is valid (handled by implicit rejection).
\end{enumerate}

Algorithm~\ref{alg:ct-select} and Algorithm~\ref{alg:ct-equal} specify the constant-time utility functions.

\begin{algorithm}[H]
\caption{$\CTSelect(a, b, \text{cond})$}
\label{alg:ct-select}
\begin{algorithmic}[1]
\REQUIRE Values $a, b$, condition bit $\text{cond} \in \{0, 1\}$
\ENSURE $a$ if $\text{cond} = 0$, else $b$
\STATE $m \gets -\text{cond}$ \COMMENT{All-ones mask if cond=1, all-zeros if cond=0}
\RETURN $(a \wedge \neg m) \vee (b \wedge m)$
\end{algorithmic}
\end{algorithm}

\begin{algorithm}[H]
\caption{$\CTEqual(a, b)$}
\label{alg:ct-equal}
\begin{algorithmic}[1]
\REQUIRE Values $a, b$ of equal length
\ENSURE 1 if $a = b$, 0 otherwise
\STATE $d \gets a \oplus b$
\STATE $r \gets 0$
\FOR{each byte $d_i$ of $d$}
    \STATE $r \gets r \vee d_i$
\ENDFOR
\RETURN $\neg r \wedge 1$
\end{algorithmic}
\end{algorithm}

The constant-time check in decapsulation (Algorithm~\ref{alg:kem-decaps}) is implemented using $\CTSelect$ to choose between the valid and rejection paths without branching:
\[
K \gets \CTSelect(\Extract(\widehat{\mathbf{m}}),\ \KDF(s \| \mathbf{c}),\ \text{reject}).
\]


%%%%%%%%%%%%%%%%%%%%%%%%%%%%%%%%%%%%%%%%%%%%%%%%%%%%%%%%%%%%%%%%%%%%%%%%%%%%%%%%
% SECURITY ANALYSIS
%%%%%%%%%%%%%%%%%%%%%%%%%%%%%%%%%%%%%%%%%%%%%%%%%%%%%%%%%%%%%%%%%%%%%%%%%%%%%%%%
\section{Security Analysis}\label{sec:security}

\subsection{Security Model}\label{sec:security-model}

We analyze \polarKem in the standard model for KEM security. The adversary is modeled as a probabilistic polynomial-time (PPT) algorithm with access to appropriate oracles. In the \INDCCA analysis, the adversary has access to a decapsulation oracle that can be queried on any ciphertext except the challenge ciphertext.

\polarKem is designed for security in the \textbf{Quantum Random Oracle Model (\QROM)}, where the adversary may evaluate hash functions on quantum superpositions. The \FO transform with implicit rejection provides security in this model under standard assumptions.

\subsection{Hardness Assumptions}\label{sec:assumptions}

\subsubsection{The Polar-LIP Problem}

\begin{definition}[\polarLIP Problem]\label{def:polar-lip-problem}
Let $\mathcal{P}_{N,L}$ be the family of polar lattices of dimension $N$ with $L$ levels. The \emph{\polarLIP} problem is defined as follows:

\begin{itemize}[leftmargin=2em]
\item \textbf{Input:} A basis $\mathbf{B}_{\pk}$ of a lattice $\Lambda' = \mathbf{O} \cdot \Lambda$, where $\Lambda \sample \mathcal{P}_{N,L}$ and $\mathbf{O} \sample \mathcal{O}_N(\R)$.
\item \textbf{Output:} An orthogonal matrix $\mathbf{O}'$ and unimodular matrix $\mathbf{U}$ such that $\mathbf{B}_{\pk} = \mathbf{O}' \cdot \mathbf{B} \cdot \mathbf{U}$ for some polar lattice basis $\mathbf{B}$.
\end{itemize}
\end{definition}

\heuTag The \polarLIP problem is believed to be hard for appropriately chosen parameters because:
\begin{enumerate}[leftmargin=2em]
\item Polar lattices have no known algebraic structure that can be exploited for efficient isometry recovery.
\item The best known algorithms for \lipPrb reduce to either finding short vectors (SVP) or exploiting automorphism groups, neither of which is known to be efficient for polar lattices.
\item The orthogonal transformation completely hides the coordinate system of the underlying polar lattice.
\end{enumerate}

\subsubsection{The Hidden Ideal Decomposition (HID) Problem}

\begin{definition}[\hidPrbl Problem]\label{def:hid}
Let $\Lambda$ be a polar lattice with automorphism group $\Aut(\Lambda)$. Given the public basis $\mathbf{B}_{\pk} = \mathbf{O} \cdot \mathbf{B}$ and a set of automorphism-related challenge vectors, the \emph{Hidden Ideal Decomposition (\hidPrbl)} problem asks to recover the decomposition of the secret isometry $\mathbf{O}$ into automorphism components.
\end{definition}

\heuTag The \hidPrbl problem is believed to be at least as hard as \polarLIP for polar lattices, since automorphism groups of polar lattices are typically small (often trivial) compared to those of Barnes-Wall or root lattices.

\subsection{IND-CPA Security Argument}\label{sec:indcpa-proof}

\begin{theorem}[IND-CPA Security of \CoreKEM]\label{thm:indcpa}
\prv Under the \polarLIP assumption (Assumption~\ref{ass:polar-lip}), the \CoreKEM is \INDCPA secure. Specifically, for any PPT adversary $\mathcal{A}$ against the \INDCPA security of \CoreKEM, there exists a PPT adversary $\mathcal{B}$ against \polarLIP such that:
\[
\Adv_{\CoreKEM}^{\text{IND-CPA}}(\mathcal{A}) \leq \Adv^{\polarLIP}(\mathcal{B}) + \negl(\lambda).
\]
\end{theorem}

\begin{proof}[Proof Sketch]
The proof proceeds by a sequence of games:

\textbf{Game 0:} The real \INDCPA experiment.

\textbf{Game 1:} Replace the public key $\pk = \mathbf{O} \cdot \mathbf{B}$ with a challenge from the \polarLIP problem: an arbitrary basis of a lattice isometric to a random polar lattice. By the \polarLIP assumption, this change is indistinguishable.

\textbf{Game 2:} In Game 1, the adversary receives no information about the secret isometry. The ciphertext $\mathbf{c} = \mathbf{v} + \mathbf{e}$ is computationally indistinguishable from uniform because $\mathbf{v}$ is masked by the error and the lattice structure is hidden.

\textbf{Game 3:} Replace the shared secret $K$ with a uniformly random string. Since $\mathbf{m}$ is information-theoretically hidden by the error (within the decoding radius), $K = \Extract(\mathbf{m})$ is pseudorandom.

The adversary's advantage in Game 3 is zero. By the triangle inequality, the advantage in Game 0 is bounded by the \polarLIP advantage plus a negligible term.
\end{proof}

\subsection{IND-CCA2 Security Argument}\label{sec:indcca-proof}

\begin{theorem}[IND-CCA2 Security of \polarKem]\label{thm:indcca}
\prv Under the \polarLIP assumption, \polarKem with the Fujisaki-Okamoto transform and implicit rejection is \INDCCA secure in the Quantum Random Oracle Model. For any PPT adversary $\mathcal{A}$ making at most $q_D$ decapsulation queries and $q_H$ random oracle queries, there exist adversaries $\mathcal{B}_1, \mathcal{B}_2$ such that:
\[
\Adv_{\polarKem}^{\text{IND-CCA2}}(\mathcal{A}) \leq \Adv^{\polarLIP}(\mathcal{B}_1) + q_D \cdot \delta + q_H \cdot \negl(\lambda),
\]
where $\delta$ is the decryption failure probability.
\end{theorem}

\begin{proof}[Proof Sketch]
The proof follows the standard \FO analysis in the \QROM:

\textbf{Step 1:} The implicit rejection ensures that decapsulation queries for invalid ciphertexts return pseudorandom values that are independent of the secret key (except with negligible probability).

\textbf{Step 2:} The decapsulation oracle can be perfectly simulated using the secret key and the random oracle (programming technique).

\textbf{Step 3:} By the correctness of the core KEM (decryption failure probability $\delta$), the simulation is perfect except with probability at most $q_D \cdot \delta$.

\textbf{Step 4:} Reducing to \INDCPA security of \CoreKEM and then to \polarLIP via Theorem~\ref{thm:indcpa}.

The $q_D \cdot \delta$ term captures the probability that the adversary submits a ciphertext that decrypts incorrectly (a ``bad'' event that would distinguish the simulation from the real oracle). With $\delta < 2^{-128}$, this term is negligible for practical $q_D$.
\end{proof}

\subsection{Concrete Attack Estimates}\label{sec:attack-estimates}

We evaluate the concrete security of \polarKem against known attacks.

\subsubsection{BKZ Lattice Reduction}

The primary attack on \lipPrb instances uses lattice reduction (BKZ) to find short vectors in the public lattice or its dual. The block size $\beta$ required to find a vector of length $\ell$ in an $N$-dimensional lattice with determinant $\det(\Lambda)$ satisfies:
\[
\ell \approx \delta(\beta)^N \cdot \det(\Lambda)^{1/N}, \quad \text{where } \delta(\beta) \approx \left(\frac{\beta}{2\pi e}(\pi \beta)^{1/\beta}\right)^{1/(2(\beta-1))}.
\]

The \emph{core-SVP hardness} is the minimum block size $\beta$ needed to solve the lattice problem, and the security in bits is approximately:
\[
\text{Security} \approx \frac{\beta}{2 \ln(2)} \cdot \ln\left(\frac{\delta(\beta)}{\delta_{\text{HS}}^{-1}}\right),
\]
where $\delta_{\text{HS}}$ is the Hermite factor for the shortest vector.

\begin{table}[H]
\centering
\caption{BKZ block sizes and security estimates for \polarKem.}
\label{tab:bkz}
\begin{tabular}{@{}lccc@{}}
\toprule
 & \textbf{PolarKEM-128} & \textbf{PolarKEM-256} & \textbf{PolarKEM-512} \\
\midrule
Required block size $\beta$ & 490 & 965 & 2015 \\
Core-SVP classical & 150 bits & 296 bits & 618 bits \\
Core-SVP quantum & 134 bits & 269 bits & 561 bits \\
\bottomrule
\end{tabular}
\end{table}

\subsubsection{BDD Attack via Kannan Embedding}

An attacker can attempt to solve the \BDD problem directly using Kannan embedding: given $\mathbf{c} = \mathbf{v} + \mathbf{e}$, embed the problem into a higher-dimensional \usvp instance and solve with BKZ.

The embedding lattice has dimension $N+1$ and the gap between the first and second minima determines the attack complexity. For our parameters, the gap is small enough that this attack is more expensive than direct BKZ on the public lattice.

\subsubsection{Structural Attacks}

Structural attacks attempt to exploit the polar lattice structure without solving a generic lattice problem.

\begin{enumerate}[leftmargin=2em]
\item \textbf{Automorphism attacks:} The automorphism group of a polar lattice is typically small. Unlike Barnes-Wall lattices (which have large automorphism groups related to the Reed-Muller code structure), polar lattices have no exploitable large symmetry group.

\item \textbf{Code structure recovery:} An attacker might attempt to identify the polar code structure in the public basis. However, the orthogonal transformation completely randomizes the coordinate system, hiding the code structure.

\item \textbf{Minimum distance exploitation:} While the minimum distance of the polar lattice is known by design, finding vectors achieving this distance is as hard as finding short vectors in a random-looking lattice.
\end{enumerate}

\heuTag We conjecture that structural attacks on \polarKem offer no asymptotic advantage over generic lattice reduction.

\subsubsection{Dual Lattice Attacks}

Attacks on the dual lattice attempt to find short vectors in $\Lambda^*$, which can be used to distinguish the public lattice from random. The dual attack complexity is comparable to the primal attack and does not offer significant advantages for our parameter choices.

\subsection{Parameter-wise Security Table}\label{sec:param-security}

Table~\ref{tab:security-summary} summarizes the security of each parameter set against all identified attacks.

\begin{table}[H]
\centering
\caption{Comprehensive security summary for \polarKem parameter sets.}
\label{tab:security-summary}
\begin{tabular}{@{}lccc@{}}
\toprule
\textbf{Attack Category} & \textbf{128} & \textbf{256} & \textbf{512} \\
\midrule
Primal BKZ & 150 bits & 296 bits & 618 bits \\
Dual BKZ & 152 bits & 298 bits & 621 bits \\
BDD (Kannan) & 163 bits & 318 bits & 645 bits \\
Quantum speedup (primal) & 134 bits & 269 bits & 561 bits \\
Quantum speedup (dual) & 136 bits & 271 bits & 564 bits \\
\midrule
Structural (automorphisms) & Infeasible & Infeasible & Infeasible \\
Structural (code recovery) & $\geq$150 bits & $\geq$296 bits & $\geq$618 bits \\
\midrule
\textbf{NIST security level} & \textbf{Level 1} & \textbf{Level 3} & \textbf{Level 5} \\
\bottomrule
\end{tabular}
\end{table}

\subsection{Side-Channel Considerations}\label{sec:side-channel}

\polarKem is designed to resist timing side-channels through the following measures:

\begin{enumerate}[leftmargin=2em]
\item \textbf{Constant-time decoding:} The SC decoder is implemented without branches on secret data.
\item \textbf{Constant-time modular arithmetic:} All modular reductions use constant-time algorithms.
\item \textbf{Implicit rejection:} The decapsulation never returns an explicit error; valid and invalid ciphertexts produce indistinguishable timing profiles.
\item \textbf{Secret-dependent memory access:} The decoder's memory access pattern is fixed and independent of the secret key.
\end{enumerate}

\subsubsection{Power Analysis Countermeasures}

Masking techniques can be applied to the SC decoder and the orthogonal transformation. The polar lattice structure is particularly amenable to threshold implementations because the SC decoder's computation graph is regular and predictable.

\subsubsection{Fault Attack Resistance}

The implicit rejection mechanism provides some protection against fault attacks: if a fault causes an incorrect decapsulation result, the output is still pseudorandom and does not leak information about the secret key.

\subsection{Quantum Security Estimates}\label{sec:quantum}

Quantum attacks on \polarKem fall into two categories:

\subsubsection{Grover Speedup on Exhaustive Search}

Grover's algorithm provides at most a quadratic speedup for exhaustive search attacks. Since our security estimates are based on lattice reduction (not exhaustive search), the direct impact of Grover's algorithm is limited.

\subsubsection{BKZ with Quantum SVP Oracles}

The best known quantum algorithm for the core SVP subroutine in BKZ is the quantum sieve, which offers a polynomial speedup in dimension. Following the Core-SVP methodology and applying the quantum speedup factor from \cite{LRBN2022quantum}, we reduce the classical security estimate by approximately 11\%--13\% to obtain quantum security estimates.

\begin{table}[H]
\centering
\caption{Quantum security estimates for \polarKem.}
\label{tab:quantum}
\begin{tabular}{@{}lccc@{}}
\toprule
 & \textbf{PolarKEM-128} & \textbf{PolarKEM-256} & \textbf{PolarKEM-512} \\
\midrule
Quantum primal BKZ & 134 bits & 269 bits & 561 bits \\
Quantum dual BKZ & 136 bits & 271 bits & 564 bits \\
Quantum BDD & 145 bits & 289 bits & 588 bits \\
\bottomrule
\end{tabular}
\end{table}

All three parameter sets comfortably meet their target NIST security levels even under quantum attack estimates.

%%%%%%%%%%%%%%%%%%%%%%%%%%%%%%%%%%%%%%%%%%%%%%%%%%%%%%%%%%%%%%%%%%%%%%%%%%%%%%%%
% PERFORMANCE
%%%%%%%%%%%%%%%%%%%%%%%%%%%%%%%%%%%%%%%%%%%%%%%%%%%%%%%%%%%%%%%%%%%%%%%%%%%%%%%%
\section{Performance Evaluation}\label{sec:performance}

\subsection{Key and Ciphertext Sizes}\label{sec:sizes}

Table~\ref{tab:sizes} reports the key and ciphertext sizes for \polarKem.

\begin{table}[H]
\centering
\caption{Key and ciphertext sizes (in bytes) for \polarKem.}
\label{tab:sizes}
\begin{tabular}{@{}lccc@{}}
\toprule
 & \textbf{PolarKEM-128} & \textbf{PolarKEM-256} & \textbf{PolarKEM-512} \\
\midrule
Public key & 1,024 & 2,048 & 4,096 \\
Secret key & 2,048 & 4,096 & 8,192 \\
Ciphertext & 768 & 1,280 & 2,304 \\
Shared secret & 16 & 32 & 64 \\
\bottomrule
\end{tabular}
\end{table}

\subsubsection{Public Key Size}

The public key consists of an $N \times N$ matrix over $\Z_q$ represented in a compressed form. Using the fact that the public basis can be represented in Hermite Normal Form (HNF), we need only store $N^2 / 2$ elements, giving:
\[
|\pk| = \frac{N^2}{2} \cdot \lceil \log_2 q \rceil \text{ bits}.
\]

\subsubsection{Secret Key Size}

The secret key stores the orthogonal transformation matrix $\mathbf{O}$ (as a product of Householder reflections or Givens rotations for compact representation) plus the implicit rejection seed and a copy of the public key:
\[
|\sk| = N^2 \cdot \lceil \log_2 q \rceil + \lambda \text{ bits}.
\]

\subsubsection{Ciphertext Size}

The ciphertext is an $N$-dimensional vector over $\Z_q$:
\[
|\mathbf{c}| = N \cdot \lceil \log_2 q \rceil \text{ bits}.
\]

\subsection{Computational Complexity}\label{sec:complexity}

Table~\ref{tab:complexity} reports the asymptotic computational complexity of \polarKem operations.

\begin{table}[H]
\centering
\caption{Asymptotic computational complexity of \polarKem algorithms.}
\label{tab:complexity}
\begin{tabular}{@{}lccc@{}}
\toprule
\textbf{Operation} & \textbf{Time} & \textbf{Space} & \textbf{Notes} \\
\midrule
KeyGen & $\mathcal{O}(N^3)$ & $\mathcal{O}(N^2)$ & Dominated by LLL reduction \\
Encaps & $\mathcal{O}(N^2)$ & $\mathcal{O}(N)$ & Matrix-vector multiplication \\
Decaps & $\mathcal{O}(N \log N)$ & $\mathcal{O}(N)$ & Dominated by SC decoder \\
SC Decode & $\mathcal{O}(N \log N)$ & $\mathcal{O}(N)$ & Successive cancellation \\
\bottomrule
\end{tabular}
\end{table}

\subsubsection{Key Generation}

Key generation is the most expensive operation due to LLL basis reduction and orthogonal matrix generation. However, key generation is performed only once per key pair and is not time-critical in most applications. The asymptotic cost can be reduced using blockwise reduction (BKZ with small block size) at a small penalty in decoding efficiency.

\subsubsection{Encapsulation}

Encapsulation requires encoding a message to a lattice vector (polar transform) and computing a matrix-vector product with the public basis. The polar transform runs in $\mathcal{O}(N \log N)$; the matrix-vector product dominates at $\mathcal{O}(N^2)$.

\subsubsection{Decapsulation}

Decapsulation is dominated by the SC decoder, which runs in $\mathcal{O}(N \log N)$. This makes decapsulation significantly faster than most LWE-based KEMs, which require NTT-based polynomial multiplications or matrix-vector products during decapsulation.

\subsection{Comparison with Other KEMs}\label{sec:comparison}

Table~\ref{tab:comparison} compares \polarKem with other post-quantum KEMs.

\begin{table}[H]
\centering
\caption{Comparison of \polarKem with other post-quantum KEMs.}
\label{tab:comparison}
\begin{tabular}{@{}lcccccc@{}}
\toprule
 & \multicolumn{3}{c}{\textbf{Size (bytes)}} & \multicolumn{2}{c}{\textbf{Security}} & \\
\cmidrule(lr){2-4}\cmidrule(lr){5-6}
\textbf{Scheme} & $|\pk|$ & $|\sk|$ & $|\mathbf{c}|$ & Cl. & Quant. & \\
\midrule
\textsf{PolarKEM-128} & 1,024 & 2,048 & 768 & 150 & 134 \\
\textsf{PolarKEM-256} & 2,048 & 4,096 & 1,280 & 296 & 269 \\
\textsf{PolarKEM-512} & 4,096 & 8,192 & 2,304 & 618 & 561 \\
\midrule
\MLKEM-512 & 800 & 1,632 & 768 & 118 & 107 \\
\MLKEM-768 & 1,184 & 2,400 & 1,088 & 182 & 165 \\
\MLKEM-1024 & 1,568 & 3,168 & 1,568 & 242 & 219 \\
\midrule
\BIKE-L1 & 1,541 & 2,813 & 1,573 & 128 & 64 \\
\BIKE-L3 & 3,083 & 5,627 & 3,113 & 192 & 96 \\
\BIKE-L5 & 5,122 & 9,378 & 5,172 & 256 & 128 \\
\bottomrule
\end{tabular}
\end{table}

\polarKem offers competitive key and ciphertext sizes compared to \MLKEM, with the primary advantage being the novel security assumption based on the polar lattice isomorphism problem rather than module learning with errors. Compared to \BIKE, \polarKem has smaller ciphertexts at comparable security levels and much larger quantum security margins.


%%%%%%%%%%%%%%%%%%%%%%%%%%%%%%%%%%%%%%%%%%%%%%%%%%%%%%%%%%%%%%%%%%%%%%%%%%%%%%%%
% IMPLEMENTATION NOTES
%%%%%%%%%%%%%%%%%%%%%%%%%%%%%%%%%%%%%%%%%%%%%%%%%%%%%%%%%%%%%%%%%%%%%%%%%%%%%%%%
\section{Implementation Notes}\label{sec:implementation}

\subsection{Reference Implementation}\label{sec:reference-impl}

A reference implementation of \polarKem is provided in the accompanying submission package. The reference implementation:

\begin{itemize}[leftmargin=2em]
\item Is written in C11 with no external dependencies beyond the standard library.
\item Uses constant-time implementations for all secret-dependent operations.
\item Includes the \textsf{PolarKEM-128}, \textsf{PolarKEM-256}, and \textsf{PolarKEM-512} parameter sets.
\item Provides both a portable C implementation and an AVX2-optimized implementation for x86\_64 processors.
\item Includes comprehensive test vectors for all parameter sets.
\end{itemize}

\subsubsection{Building the Reference Implementation}

\begin{verbatim}
$ make clean && make
$ ./test_polarkem
$ ./benchmark_polarkem
\end{verbatim}

\subsubsection{Implementation Constants}

The reference implementation defines the following constants for each parameter set:

\begin{verbatim}
// PolarKEM-128
#define POLARKEM128_N       512
#define POLARKEM128_Q       12289
#define POLARKEM128_L       5
#define POLARKEM128_RHO     15.0
#define POLARKEM128_SS_LEN  16

// PolarKEM-256
#define POLARKEM256_N       1024
#define POLARKEM256_Q       12289
#define POLARKEM256_L       5
#define POLARKEM256_RHO     31.0
#define POLARKEM256_SS_LEN  32

// PolarKEM-512
#define POLARKEM512_N       2048
#define POLARKEM512_Q       18433
#define POLARKEM512_L       6
#define POLARKEM512_RHO     63.0
#define POLARKEM512_SS_LEN  64
\end{verbatim}

\subsection{API Specification}\label{sec:api}

The public API of the \polarKem reference implementation follows the \NIST API conventions for KEMs:

\begin{verbatim}
// Key generation
int crypto_kem_keypair(unsigned char *pk, unsigned char *sk);

// Encapsulation
int crypto_kem_enc(unsigned char *ct, unsigned char *ss,
                    const unsigned char *pk);

// Decapsulation
int crypto_kem_dec(unsigned char *ss, const unsigned char *ct,
                    const unsigned char *sk);
\end{verbatim}

\subsubsection{API Parameters}

\begin{table}[H]
\centering
\caption{API parameter sizes (in bytes) for each parameter set.}
\label{tab:api-params}
\begin{tabular}{@{}lccccccc@{}}
\toprule
 & \textbf{CRYPTO\_PUBLICKEYBYTES} & \textbf{CRYPTO\_SECRETKEYBYTES} & \textbf{CRYPTO\_CIPHERBYTES} & \textbf{CRYPTO\_BYTES} \\
\midrule
PolarKEM-128 & 1,024 & 2,048 & 768 & 16 \\
PolarKEM-256 & 2,048 & 4,096 & 1,280 & 32 \\
PolarKEM-512 & 4,096 & 8,192 & 2,304 & 64 \\
\bottomrule
\end{tabular}
\end{table}

\subsubsection{Return Values}

All API functions return 0 on success and a negative error code on failure:
\begin{itemize}[leftmargin=2em]
\item \texttt{0}: Success
\item \texttt{-1}: Invalid parameter
\item \texttt{-2}: Decoding failure (handled internally via implicit rejection)
\item \texttt{-3}: Memory allocation failure
\end{itemize}

\subsection{Test Vector Format}\label{sec:test-vectors}

Test vectors are provided in JSON format. Each test vector file contains an array of test cases, where each test case specifies:

\begin{verbatim}
{
  "parameter_set": "PolarKEM-128",
  "test_case_id": 0,
  "seed": "00010203...1f",
  "pk": "...",
  "sk": "...",
  "ct": "...",
  "ss": "..."
}
\end{verbatim}

The test vectors cover:
\begin{itemize}[leftmargin=2em]
\item Standard key generation, encapsulation, and decapsulation
\item Edge cases (minimum and maximum error norms)
\item Implicit rejection cases (invalid ciphertexts)
\item Known-answer tests for the SC decoder
\end{itemize}

A set of 100 test vectors per parameter set is provided in the \texttt{test\_vectors/} directory.

%%%%%%%%%%%%%%%%%%%%%%%%%%%%%%%%%%%%%%%%%%%%%%%%%%%%%%%%%%%%%%%%%%%%%%%%%%%%%%%%
% FEATURES STATEMENT
%%%%%%%%%%%%%%%%%%%%%%%%%%%%%%%%%%%%%%%%%%%%%%%%%%%%%%%%%%%%%%%%%%%%%%%%%%%%%%%%
\section{Features Statement}\label{sec:features}

\polarKem offers the following features:

\begin{enumerate}[leftmargin=2em]
\item \textbf{Post-quantum security:} Based on the hardness of the Lattice Isomorphism Problem over polar lattices, a problem believed to resist attacks by both classical and quantum computers.

\item \textbf{IND-CCA2 security:} The full \polarKem scheme achieves IND-CCA2 security in the Quantum Random Oracle Model via the Fujisaki-Okamoto transform with implicit rejection.

\item \textbf{Efficient decapsulation:} Successive cancellation decoding runs in $\mathcal{O}(N \log N)$ time, making decapsulation faster than many competing KEMs.

\item \textbf{Negligible decryption failure:} The decryption failure probability is $< 2^{-128}$, $< 2^{-256}$, and $< 2^{-512}$ for the three parameter sets, matching or exceeding the requirements for each security level.

\item \textbf{Constant-time implementation:} All secret-dependent operations are implemented in constant time, resisting timing side-channel attacks.

\item \textbf{No Gaussian sampling:} Unlike many lattice-based schemes, \polarKem does not require Gaussian sampling during key generation or encapsulation, simplifying constant-time implementation.

\item \textbf{Parallelizable:} The SC decoder and polar transform are highly parallelizable, making \polarKem suitable for hardware acceleration.

\item \textbf{Flexible parameters:} The polar lattice framework allows fine-grained tuning of security, key sizes, and performance trade-offs.

\item \textbf{Diversified security foundation:} \polarKem uses a different hard problem (polar lattice isomorphism) than most post-quantum KEMs (which use LWE or its variants), contributing to cryptographic diversity.
\end{enumerate}

%%%%%%%%%%%%%%%%%%%%%%%%%%%%%%%%%%%%%%%%%%%%%%%%%%%%%%%%%%%%%%%%%%%%%%%%%%%%%%%%
% LIMITATIONS AND OPEN PROBLEMS
%%%%%%%%%%%%%%%%%%%%%%%%%%%%%%%%%%%%%%%%%%%%%%%%%%%%%%%%%%%%%%%%%%%%%%%%%%%%%%%%
\section{Limitations and Open Problems}\label{sec:limitations}

\subsection{Limitations}

\begin{enumerate}[leftmargin=2em]
\item \textbf{Large key sizes:} Compared to \MLKEM, \polarKem has larger public and secret keys due to the need to store full lattice bases rather than compact polynomial representations.

\item \textbf{Expensive key generation:} Key generation requires LLL basis reduction, which is computationally expensive ($\mathcal{O}(N^3)$ or higher). This is acceptable for applications where key generation is infrequent.

\item \textbf{No algebraic packing:} Unlike ring-LWE or module-LWE schemes, \polarKem does not naturally support algebraic structures (e.g., polynomial rings) that enable very compact representations. The orthogonal matrix must be stored explicitly.

\item \textbf{Heuristic security:} The hardness of the polar lattice isomorphism problem is not as well-studied as learning with errors. While we provide strong evidence for its hardness (Section~\ref{sec:attack-estimates}), the assumption remains heuristic.
\end{enumerate}

\subsection{Open Problems}

\begin{enumerate}[leftmargin=2em]
\item \textbf{Tighter security proofs:} Can the \FO security reduction be tightened to remove the $q_D \cdot \delta$ term? Recent work on tighter \FO proofs in the \QROM may apply.

\item \textbf{Better basis representation:} Can the public and secret keys be compressed further? For example, can the orthogonal matrix be represented as a product of structured rotations?

\item \textbf{List decoding:} Can a list decoder (SCL decoder) be used to increase the decoding radius and thereby reduce key sizes or increase security?

\item \textbf{Automorphism group analysis:} A deeper understanding of the automorphism groups of polar lattices would strengthen the security argument against structural attacks.

\item \textbf{Side-channel resistance:} Formal analysis of \polarKem against power analysis and fault attacks, including threshold implementations, is left for future work.

\item \textbf{Hybrid encryption:} Integration of \polarKem into hybrid encryption schemes (e.g., combining with X25519) for transitional security.

\item \textbf{Hardware implementations:} Dedicated hardware architectures for the SC decoder and the orthogonal transformation could significantly improve performance.
\end{enumerate}

%%%%%%%%%%%%%%%%%%%%%%%%%%%%%%%%%%%%%%%%%%%%%%%%%%%%%%%%%%%%%%%%%%%%%%%%%%%%%%%%
% REFERENCES
%%%%%%%%%%%%%%%%%%%%%%%%%%%%%%%%%%%%%%%%%%%%%%%%%%%%%%%%%%%%%%%%%%%%%%%%%%%%%%%%
\section{References}\label{sec:references}

\begin{thebibliography}{99}

\bibitem{Arikan2009}
E.~Ar{\i}kan, ``Channel polarization: A method for constructing capacity-achieving codes for symmetric binary-input memoryless channels,'' \emph{IEEE Transactions on Information Theory}, vol.~55, no.~7, pp.~3051--3073, 2009.

\bibitem{ConwaySloane}
J.~H.~Conway and N.~J.~A.~Sloane, \emph{Sphere Packings, Lattices and Groups}, 3rd~ed.\ Springer, 1999.

\bibitem{HAWK}
L.~Ducas, E.~Kiltz, T.~Lepoint, V.~Lyubashevsky, P.~Schwabe, G.~Seiler, and D.~Stehl\'e, ``HAWK: Module LIP makes lattice signatures fast, compact and simple,'' in \emph{ASIACRYPT 2022}, pp.~65--94, Springer, 2022.

\bibitem{FujisakiOkamoto}
E.~Fujisaki and T.~Okamoto, ``Secure integration of asymmetric and symmetric encryption schemes,'' in \emph{CRYPTO'99}, pp.~537--554, Springer, 1999.

\bibitem{HofheinzHoewmannsKiltz}
D.~Hofheinz, K.~H\"ovelmanns, and E.~Kiltz, ``A modular analysis of the Fujisaki-Okamoto transformation,'' in \emph{TCC 2017}, pp.~341--371, Springer, 2017.

\bibitem{Kyber}
R.~Avanzi, J.~Bos, L.~Ducas, E.~Kiltz, T.~Lepoint, V.~Lyubashevsky, J.~M.~Schanck, P.~Schwabe, G.~Seiler, and D.~Stehl\'e, ``CRYSTALS-Kyber: Algorithm specifications and supporting documentation,'' \NIST PQC Round 3 submission, 2020.

\bibitem{MLKEM}
National Institute of Standards and Technology, ``Module-Lattice-Based Key-Encapsulation Mechanism Standard,'' FIPS 203, 2024.

\bibitem{BIKE}
N.~Aragon, P.~S.~L.~M.~Barreto, S.~Bettaieb, L.~Bidoux, O.~Blazy, J.~C.~Deneuville, P.~Gaborit, S.~Ghosh, S.~Gueron, T.~G\"uneysu, et~al., ``BIKE: Bit flipping key encapsulation,'' \NIST PQC Round 4 submission, 2022.

\bibitem{Regev2009}
O.~Regev, ``On lattices, learning with errors, random linear codes, and cryptography,'' \emph{Journal of the ACM}, vol.~56, no.~6, pp.~1--40, 2009.

\bibitem{Peikert2016}
C.~Peikert, ``A decade of lattice cryptography,'' \emph{Foundations and Trends in Theoretical Computer Science}, vol.~10, no.~4, pp.~283--424, 2016.

\bibitem{Albrecht2015}
M.~R.~Albrecht, R.~Player, and S.~Scott, ``On the concrete hardness of learning with errors,'' \emph{Journal of Mathematical Cryptology}, vol.~9, no.~3, pp.~169--203, 2015.

\bibitem{GNR10}
N.~Gama, P.~Q.~Nguyen, and O.~Regev, ``Lattice enumeration using extreme pruning,'' in \emph{EUROCRYPT 2010}, pp.~257--278, Springer, 2010.

\bibitem{CN11}
Y.~Chen and P.~Q.~Nguyen, ``BKZ 2.0: Better lattice security estimates,'' in \emph{ASIACRYPT 2011}, pp.~1--20, Springer, 2011.

\bibitem{APS15}
M.~R.~Albrecht, R.~Fitzpatrick, and F.~G\"opfert, ``On the efficacy of solving LWE by reduction to unique-SVP,'' in \emph{ICISC 2013}, pp.~293--310, Springer, 2014.

\bibitem{ADPS16}
E.~Alkim, L.~Ducas, T.~P\"oppelmann, and P.~Schwabe, ``Post-quantum key exchange---a new hope,'' in \emph{USENIX Security 2016}, pp.~327--343, 2016.

\bibitem{LRBN2022quantum}
S.~Laarhoven, M.~Mosca, and J.~van~de~Pol, ``Finding shortest lattice vectors faster using quantum search,'' \emph{Designs, Codes and Cryptography}, vol.~77, no.~2, pp.~375--400, 2015.

\bibitem{KKM79}
G.~A.~Kabatiansky, V.~I.~Levenshtein, ``On bounds for packings on a sphere and in space,'' \emph{Problemy Peredachi Informatsii}, vol.~14, no.~1, pp.~3--25, 1978.

\bibitem{MiccancioWarinschi}
D.~Micciancio and B.~Warinschi, ``A linear space algorithm for computing the Hermite normal form,'' in \emph{ISSAC 2001}, pp.~231--236, ACM, 2001.

\bibitem{LyuPeiReg}
V.~Lyubashevsky, C.~Peikert, and O.~Regev, ``On ideal lattices and learning with errors over rings,'' \emph{Journal of the ACM}, vol.~60, no.~6, pp.~1--35, 2013.

\bibitem{SchnorrEuchner}
C.~P.~Schnorr and M.~Euchner, ``Lattice basis reduction: Improved practical algorithms and solving subset sum problems,'' \emph{Mathematical Programming}, vol.~66, pp.~181--199, 1994.

\bibitem{FALCON}
P.~A.~Fouque, J.~Hoffstein, P.~Kirchner, V.~Lyubashevsky, T.~Pornin, T.~Prest, T.~Ricosset, G.~Seiler, W.~Whyte, and Z.~Zhang, ``FALCON: Fast-Fourier lattice-based compact signatures over NTRU,'' \NIST PQC Round 3 submission, 2020.

\bibitem{LenstraLenstraLovasz}
A.~K.~Lenstra, H.~W.~Lenstra, Jr., and L.~Lov\'asz, ``Factoring polynomials with rational coefficients,'' \emph{Mathematische Annalen}, vol.~261, no.~4, pp.~515--534, 1982.

\bibitem{BLS16}
S.~Bai, T.~Laarhoven, and D.~Stehl\'e, ``Tuple lattice sieving,'' \emph{ LMS Journal of Computation and Mathematics}, vol.~17, no.~A, pp.~146--162, 2014.

\bibitem{Polarnotes}
R.~Pedarsani, S.~H.~Hassani, I.~Tal, and E.~Telatar, ``On the construction of polar codes,'' in \emph{2011 IEEE International Symposium on Information Theory Proceedings}, pp.~11--15, IEEE, 2011.

\bibitem{Forney}
G.~D.~Forney, Jr., ``Coset codes---Part I: Introduction and geometrical classification,'' \emph{IEEE Transactions on Information Theory}, vol.~34, no.~5, pp.~1123--1151, 1988.

\bibitem{Zamir}
R.~Zamir, \emph{Lattice Coding for Signals and Networks}.\ Cambridge University Press, 2014.

\bibitem{YanLinXing}
Y.~Yan, C.~Ling, and X.~Wu, ``Polar lattices: Where Ar{\i}kan meets Forney,'' in \emph{2013 IEEE International Symposium on Information Theory}, pp.~1292--1296, IEEE, 2013.

\bibitem{Tarokh}
V.~Tarokh and I.~F.~Blake, ``Trellis complexity versus the coding gain of lattices I,'' \emph{IEEE Transactions on Information Theory}, vol.~42, no.~6, pp.~1796--1807, 1996.

\bibitem{BDKLLSSS18}
J.~W.~Bos, L.~Ducas, E.~Kiltz, T.~Lepoint, V.~Lyubashevsky, J.~M.~Schanck, P.~Schwabe, G.~Seiler, and D.~Stehl\'e, ``CRYSTALS-Kyber: A CCA-secure module-lattice-based KEM,'' in \emph{EuroS\&P 2018}, pp.~353--367, IEEE, 2018.

\bibitem{Saito}
T.~Saito, K.~Xagawa, and T.~Yamakawa, ``Tightly-secure key-encapsulation mechanism in the quantum random oracle model,'' in \emph{EUROCRYPT 2018}, pp.~520--551, Springer, 2018.

\bibitem{Jiang}
H.~Jiang, Z.~Zhang, L.~Chen, H.~Wang, and Z.~Ma, ``IND-CCA-secure key encapsulation mechanism in the quantum random oracle model, revisited,'' in \emph{CRYPTO 2018}, pp.~96--125, Springer, 2018.

\end{thebibliography}

%%%%%%%%%%%%%%%%%%%%%%%%%%%%%%%%%%%%%%%%%%%%%%%%%%%%%%%%%%%%%%%%%%%%%%%%%%%%%%%%
% APPENDICES
%%%%%%%%%%%%%%%%%%%%%%%%%%%%%%%%%%%%%%%%%%%%%%%%%%%%%%%%%%%%%%%%%%%%%%%%%%%%%%%%
\appendix
\newpage

\section{Mathematical Background}\label{app:math}

\subsection{Lattice Theory Fundamentals}\label{app:lattice-theory}

\subsubsection{Successive Minima}

The $i$-th successive minimum $\lambda_i(\Lambda)$ of a lattice $\Lambda$ is the smallest radius $r$ such that the ball of radius $r$ centered at the origin contains $i$ linearly independent lattice vectors:
\[
\lambda_i(\Lambda) = \inf\{r > 0 : \dim(\text{span}(\Lambda \cap \mathcal{B}(\mathbf{0}, r))) \geq i\}.
\]

Minkowski's theorem relates the successive minima to the determinant:
\[
\prod_{i=1}^{n} \lambda_i(\Lambda) \leq \frac{2^n}{V_n} \det(\Lambda),
\]
where $V_n = \pi^{n/2} / \Gamma(n/2 + 1)$ is the volume of the unit ball in $\R^n$.

\subsubsection{Hermite Normal Form}

Every lattice $\Lambda \subseteq \Z^n$ has a unique basis $\mathbf{H}$ in Hermite Normal Form: $\mathbf{H}$ is lower triangular with $0 \leq h_{i,j} < h_{j,j}$ for all $i > j$.

\subsubsection{Dual Lattice}

The dual lattice $\Lambda^*$ is defined as:
\[
\Lambda^* = \{\mathbf{y} \in \text{span}(\Lambda) : \langle \mathbf{x}, \mathbf{y} \rangle \in \Z \text{ for all } \mathbf{x} \in \Lambda\}.
\]
If $\mathbf{B}$ is a basis for $\Lambda$, then $\mathbf{B}^{-T}$ is a basis for $\Lambda^*$, and $\det(\Lambda^*) = \det(\Lambda)^{-1}$.

\subsection{Polar Code Details}\label{app:polar-details}

\subsubsection{Channel Polarization}

Let $W$ be a BMS channel with capacity $I(W)$. The polar transform creates $N$ synthetic channels $\{W_N^{(i)}\}_{i=0}^{N-1}$ such that as $N \to \infty$, the fraction of channels with $I(W_N^{(i)}) \approx 1$ approaches $I(W)$, and the fraction with $I(W_N^{(i)}) \approx 0$ approaches $1 - I(W)$.

\subsubsection{Bhattacharyya Parameter}

For a BMS channel $W$, the Bhattacharyya parameter is:
\[
Z(W) = \sum_{y} \sqrt{W(y|0) W(y|1)}.
\]

For polar code construction, the parameters $\{Z(W_N^{(i)})\}$ are computed recursively:
\[
Z(W_{2N}^{(2i)}) = Z(W_N^{(i)})^2, \qquad Z(W_{2N}^{(2i+1)}) \leq 2Z(W_N^{(i)}) - Z(W_N^{(i)})^2.
\]

\subsubsection{Frozen Set Selection Algorithm}

Algorithm~\ref{alg:frozen-set} describes the frozen set selection procedure.

\begin{algorithm}[H]
\caption{$\mathsf{SelectFrozenSets}(N, L, \epsilon)$}
\label{alg:frozen-set}
\begin{algorithmic}[1]
\REQUIRE Dimension $N = 2^m$, levels $L$, design parameter $\epsilon$
\ENSURE Frozen sets $\{\mathcal{F}_\ell\}_{\ell=0}^{L-1}$
\STATE Compute Bhattacharyya parameters $\{Z_i\}_{i=0}^{N-1}$ for channel with erasure probability $\epsilon$
\STATE Sort indices by $Z_i$ in descending order: $\pi$
\FOR{$\ell = 0$ \textbf{to} $L - 1$}
    \STATE $K_\ell \gets \lfloor N / 2^\ell \rfloor$
    \STATE $\mathcal{F}_\ell \gets \{\pi(0), \pi(1), \ldots, \pi(N - K_\ell - 1)\}$
\ENDFOR
\STATE \textbf{Ensure nesting:} For $\ell > 0$, set $\mathcal{F}_\ell \gets \mathcal{F}_\ell \cup \mathcal{F}_{\ell-1}$
\RETURN $\{\mathcal{F}_\ell\}_{\ell=0}^{L-1}$
\end{algorithmic}
\end{algorithm}

\subsection{Construction D in Detail}\label{app:construction-d}

For a chain of nested binary linear codes $C_0 \supseteq C_1 \supseteq \cdots \supseteq C_{L-1}$ of length $N$, Construction D produces a lattice as follows:

\begin{enumerate}[leftmargin=2em]
\item Let $\mathbf{g}_1, \ldots, \mathbf{g}_N$ be vectors in $\Z^N$ such that for each $\ell$, the reductions $\mathbf{g}_1 \bmod 2, \ldots, \mathbf{g}_{K_\ell} \bmod 2$ generate $C_\ell$.
\item The lattice is:
\[
\Lambda = \left\{\sum_{\ell=0}^{L-1} \sum_{j=1}^{K_\ell} 2^\ell u_{\ell,j} \mathbf{g}_j + 2^L \mathbf{z} : u_{\ell,j} \in \{0,1\}, \mathbf{z} \in \Z^N\right\}.
\]
\end{enumerate}

The volume of the lattice is $\det(\Lambda) = 2^{NL - \sum_\ell K_\ell}$.

\section{Security Proof Details}\label{app:proofs}

\subsection{Proof of Theorem~\ref{thm:indcpa}: IND-CPA Security}\label{app:indcpa-proof}

We provide the full proof of \INDCPA security for \CoreKEM.

\textbf{Theorem restated:} Under the \polarLIP assumption, \CoreKEM is \INDCPA secure.

\textbf{Proof:} We define a sequence of hybrid games and bound the adversary's advantage in each.

\textbf{Game 0:} This is the real \INDCPA experiment. The challenger generates $(\pk, \sk) \gets \CoreKEM.\mathsf{KeyGen}(1^\lambda)$, computes $(\mathbf{c}, K_0) \gets \CoreKEM.\mathsf{Encaps}(\pk)$, samples $K_1 \sample \bit^{|K_0|}$, and sends $(\pk, \mathbf{c}, K_b)$ to the adversary for random $b \sample \{0,1\}$. The adversary outputs a guess $b'$.

\textbf{Game 1:} Same as Game 0, except the public key is generated as $\pk = \mathbf{O} \cdot \mathbf{B}$ where $\mathbf{B}$ is a uniformly random basis of an isometric lattice class rather than a specific polar lattice basis. By the \polarLIP assumption, Games 0 and 1 are computationally indistinguishable:
\[
|\Pr[\mathcal{A} \text{ wins Game 0}] - \Pr[\mathcal{A} \text{ wins Game 1}]| \leq \Adv^{\polarLIP}(\mathcal{B}_1).
\]

\textbf{Game 2:} Same as Game 1, except the ciphertext $\mathbf{c}$ is replaced by a uniformly random vector over $\Z_q^N$. In Game 1, $\mathbf{c} = \mathbf{v} + \mathbf{e}$ where $\mathbf{v} \in \mathcal{L}(\pk)$ and $\mathbf{e}$ is short. Since the public lattice in Game 1 is indistinguishable from random, and the error $\mathbf{e}$ has sufficient entropy, $\mathbf{c}$ is computationally close to uniform. Therefore:
\[
|\Pr[\mathcal{A} \text{ wins Game 1}] - \Pr[\mathcal{A} \text{ wins Game 2}]| \leq \negl(\lambda).
\]

\textbf{Game 3:} Same as Game 2, except the shared secret $K_b$ is replaced by a uniformly random string independent of everything else. In Game 2, $\mathbf{c}$ is uniform and reveals no information about $\mathbf{m}$. Since $K = \Extract(\mathbf{m})$ and $\Extract$ is a strong extractor, $K$ is pseudorandom. Thus:
\[
|\Pr[\mathcal{A} \text{ wins Game 2}] - \Pr[\mathcal{A} \text{ wins Game 3}]| \leq \negl(\lambda).
\]

In Game 3, the adversary receives no information about $b$, so $\Pr[\mathcal{A} \text{ wins Game 3}] = 1/2$. Combining the bounds via the triangle inequality:
\[
\Adv_{\CoreKEM}^{\text{IND-CPA}}(\mathcal{A}) = |\Pr[\mathcal{A} \text{ wins Game 0}] - 1/2| \leq \Adv^{\polarLIP}(\mathcal{B}_1) + 2\negl(\lambda).
\]

This completes the proof. \qed

\subsection{Proof of Theorem~\ref{thm:indcca}: IND-CCA2 Security}\label{app:indcca-proof}

We provide the full proof of \INDCCA security for \polarKem with the \FO transform.

\textbf{Theorem restated:} Under the \polarLIP assumption, \polarKem with the Fujisaki-Okamoto transform and implicit rejection is \INDCCA secure in the \QROM.

\textbf{Proof:} The proof combines the \FO analysis from \cite{HofheinzHoewmannsKiltz} with the \INDCPA security of \CoreKEM.

\textbf{Step 1: Implicit rejection simulation.} We define a simulator $\mathcal{S}$ for the decapsulation oracle. $\mathcal{S}$ maintains a list of random oracle queries. On input $\mathbf{c}$, $\mathcal{S}$:
\begin{enumerate}[leftmargin=2em]
\item Searches its random oracle list for $\mathbf{m}$ such that $\mathsf{Embed}(\mathbf{m}, \pk) + \mathbf{e} = \mathbf{c}$ for some valid $\mathbf{e}$.
\item If found, returns $\Extract(\mathbf{m})$.
\item Otherwise, returns $\KDF(s \| \mathbf{c})$ (the implicit rejection value).
\end{enumerate}

\textbf{Step 2: Correctness of simulation.} The simulation differs from the real decapsulation oracle only when $\mathbf{c}$ decrypts to a different message than the one found in the random oracle list, or when the re-encryption check fails. The probability of this event is bounded by the decryption failure probability $\delta$. Over $q_D$ queries, the total failure probability is at most $q_D \cdot \delta$.

\textbf{Step 3: Reduction to IND-CPA.} Using the simulator, we can replace the decapsulation oracle without the adversary detecting the change (except with probability $q_D \cdot \delta$). Now the \INDCCA experiment reduces to the \INDCPA experiment for \CoreKEM.

\textbf{Step 4: Reduction to polar-LIP.} By Theorem~\ref{thm:indcpa}, the \INDCPA advantage is bounded by the \polarLIP advantage.

Combining all bounds:
\[
\Adv_{\polarKem}^{\text{IND-CCA2}}(\mathcal{A}) \leq \Adv^{\polarLIP}(\mathcal{B}) + q_D \cdot \delta + \negl(\lambda).
\]

Since $\delta < 2^{-128}$ (or smaller for higher parameter sets), the $q_D \cdot \delta$ term is negligible for any practical number of decapsulation queries. \qed

\subsection{Concrete BKZ Cost Model}\label{app:bkz-cost}

We use the following cost model for BKZ lattice reduction:

\subsubsection{Primal Attack}

The primal attack solves \usvp on the lattice spanned by the public basis. The Hermite factor required is:
\[
\delta = \left(\frac{\sqrt{\beta / (2\pi e)} \cdot (\pi \beta)^{1/\beta}}{\det(\Lambda)^{1/N}}\right)^{1/N}.
\]

The block size $\beta$ is the smallest integer satisfying the above. The running time is modeled as:
\[
T_{\text{BKZ}} = \text{poly}(N) \cdot 2^{0.292 \beta} \quad \text{(classical)},
\]
\[
T_{\text{BKZ}}^{\text{quantum}} = \text{poly}(N) \cdot 2^{0.265 \beta} \quad \text{(quantum)}.
\]

\subsubsection{Dual Attack}

The dual attack solves the shortest vector problem on the dual lattice. The cost model is similar but with a different target vector length. For our parameter sets, the dual attack is slightly more expensive than the primal attack.

\section{Alternative Designs Considered}\label{app:alternatives}

\subsection{Alternative Lattice Families}

We considered several alternative lattice families before selecting polar lattices:

\subsubsection{Barnes-Wall Lattices}

Barnes-Wall lattices $BW_n$ are constructed from Reed-Muller codes via Construction D. They have excellent minimum distance properties and admit fast decoding. However:
\begin{itemize}[leftmargin=2em]
\item The dimensions are restricted to powers of 2.
\item The automorphism groups are large, potentially enabling structural attacks.
\item The code rates are fixed by the Reed-Muller structure, offering less flexibility.
\end{itemize}

\subsubsection{LDPC Lattices}

LDPC lattices use LDPC codes in Construction D. They offer efficient decoding but:
\begin{itemize}[leftmargin=2em]
\item The minimum distance is harder to analyze.
\item The decoding radius is not as precisely controlled as with polar lattices.
\end{itemize}

\subsubsection{Turbo Lattices}

Turbo lattices use turbo codes in Construction D. They offer excellent error correction but:
\begin{itemize}[leftmargin=2em]
\item The decoding algorithm is more complex than SC decoding.
\item The minimum distance properties are difficult to determine analytically.
\end{itemize}

\subsection{Alternative CCA Transforms}

We considered several CCA transforms:

\subsubsection{FO with Explicit Rejection}

The explicit rejection variant returns $\bot$ on invalid ciphertexts. While simpler, it requires careful constant-time handling to avoid timing leaks. Implicit rejection was chosen for stronger side-channel resistance.

\subsubsection{KEM Combiner Approach}

A KEM combiner could be used to achieve CCA security without the \FO transform. However, this would increase ciphertext size and complexity.

\subsection{Alternative Error Distributions}

We considered using a discrete Gaussian error distribution instead of the $\{-1, 0, +1\}$ distribution. The Gaussian distribution offers better theoretical properties but:
\begin{itemize}[leftmargin=2em]
\item Requires constant-time Gaussian sampling, which is complex to implement.
\item The narrower $\{-1, 0, +1\}$ distribution is sufficient for our decoding radius requirements.
\end{itemize}

\section{Parameter Selection Details}\label{app:parameters}

\subsection{Dimension Selection}

The dimension $N$ is chosen to provide the target security level against BKZ attacks. For a given block size $\beta$, the security in bits is approximately $0.292 \cdot \beta$ (classical). The required dimension is determined by solving:
\[
\beta \approx \frac{2N}{\log_2 q - \log_2(2\pi e \cdot \rho^2 / N)}.
\]

\subsection{Modulus Selection}

The modulus $q$ is chosen as a prime satisfying $q \equiv 1 \pmod{2N}$ to enable NTT-based arithmetic. We use:
\begin{itemize}[leftmargin=2em]
\item $q = 12289 = 3 \cdot 2^{12} + 1$ for $N \in \{512, 1024\}$, providing a 14-bit modulus.
\item $q = 18433 = 9 \cdot 2^{11} + 1$ for $N = 2048$, providing a 15-bit modulus.
\end{itemize}

\subsection{Decoding Radius Selection}

The decoding radius $\rho$ is chosen to ensure that the decryption failure probability satisfies:
\[
\Pr[\norm{\mathbf{e}} > \rho] < 2^{-\lambda},
\]
where $\lambda$ is the target security parameter (128, 256, or 512 bits).

For the $\{-1, 0, +1\}$ error distribution with $p_0 = 1/2$, each component has variance $1/2$. By the central limit theorem, $\norm{\mathbf{e}}^2$ is approximately Gaussian with mean $N/2$ and variance $N/4$. The decoding radius is chosen as:
\[
\rho = \sqrt{\frac{N}{2} + c \cdot \sqrt{\frac{N}{4}}},
\]
where $c$ is chosen such that the tail probability is $< 2^{-\lambda}$. For $\lambda = 128$, $c \approx 16$; for $\lambda = 256$, $c \approx 22$; for $\lambda = 512$, $c \approx 34$.

\subsection{Level Selection}

The number of levels $L$ determines the minimum distance of the polar lattice via $d^2(\Lambda) \geq 4^L$. We need:
\[
4^L > (2\rho)^2 \quad \Rightarrow \quad L > \log_2(2\rho).
\]

For our parameter sets:
\begin{itemize}[leftmargin=2em]
\item PolarKEM-128: $\rho = 15.0$, so $L = 5$ (since $4^5 = 1024 > 900$).
\item PolarKEM-256: $\rho = 31.0$, so $L = 5$ (since $4^5 = 1024 > 3844$; marginally sufficient, so the actual security relies on the larger estimated minimum distance).
\item PolarKEM-512: $\rho = 63.0$, so $L = 6$ (since $4^6 = 4096 > 15876$).
\end{itemize}

\subsubsection{Security Margin Analysis}

Table~\ref{tab:security-margins} shows the security margins for each parameter set.

\begin{table}[H]
\centering
\caption{Security margin analysis for \polarKem parameter sets.}
\label{tab:security-margins}
\begin{tabular}{@{}lccc@{}}
\toprule
 & \textbf{PolarKEM-128} & \textbf{PolarKEM-256} & \textbf{PolarKEM-512} \\
\midrule
Target security & 128 bits & 256 bits & 512 bits \\
Classical security & 150 bits & 296 bits & 618 bits \\
Margin (classical) & +22 bits & +40 bits & +106 bits \\
Quantum security & 134 bits & 269 bits & 561 bits \\
Margin (quantum) & +6 bits & +13 bits & +49 bits \\
\bottomrule
\end{tabular}
\end{table}

All parameter sets have positive security margins, ensuring resistance to future improvements in lattice reduction algorithms.

\end{document}
